
\clearpage
% \appendix
 
\begin{appendices}

% \begin{table*}[t]
% \centering
% \small
% \caption{Appendix roadmap.}
% \label{tab:appendix-roadmap}
% \begin{tabular}{p{0.13\linewidth} p{0.78\linewidth}}
% \toprule
% Appendix & What it covers \\
% \midrule
% A & Out-of-core evaluation: large circuits, spill, and sparsity effects. \\
% B--C & Routing ablations and training/inference overhead. \\
% D & Tensor-aware DBMS discussion and preliminary SciDB results. \\
% E & Support for SupermarQ, MQT Bench, and QASM Bench. \\
% F & SQLite vs. Qiskit Aer profiling for memory and runtime. \\
% G--H & DBMS tuning, indexing, partitioning, and CTE materialization. \\
% I & Reproducibility details for circuits, spill metrics, and tuning. \\
% J--K & Generated-circuit characterization and subcircuit-generation details. \\
% L & Extra SQL-only feature results and preprocessing scalability. \\
% M--N & Implementation details, Qiskit Aer selector, and feature schema. \\
% \bottomrule
% \end{tabular}
% \end{table*}
\appendix

\noindent Table~\ref{tab:appendix-roadmap} provides a compact roadmap of the appendix material.

\begin{table}[H]
\centering
\caption{Appendix roadmap.}
\label{tab:appendix-roadmap}
\renewcommand{\arraystretch}{1.08}
\begin{tabularx}{\columnwidth}{@{}p{0.17\columnwidth}X@{}}
\toprule
Appendix & Main content \\
\midrule
A & Out-of-core evaluation and spill behavior. \\
B--C & Routing ablations and training/inference overhead. \\
D--E & Tensor-aware DBMS discussion and support for external benchmark suites. \\
F--H & SQLite--Aer profiling, DBMS-specific tuning, and physical-design discussion. \\
I--K & Reproducibility details, generated-workload characterization, and subcircuit generation. \\
L--O & SQL-only results, preprocessing scalability, implementation details, Qiskit Aer selector, and feature schema. \\
\bottomrule
\end{tabularx}
\end{table}

 
 
% We evaluate dense QFT circuits from 3 to 26 qubits under split SQL execution on PostgreSQL, DuckDB, and SQLite, using memory limits of 4, 8, and 16~GB. We have chosen QFT as it is intentionally challenging: it quickly produces dense intermediate states and therefore allows us to inspect whether an RDBMS can complete the simulation by spilling intermediate results to disk. All configurations complete except one: DuckDB at the largest circuit (26 qubits) under the smallest 4~GB memory limit terminates with an out-of-memory (OOM) error after 51 gates. 

 

% \medskip
% \noindent\textbf{OOC parameters.}
% For this OOC sweep, we generate one QFT circuit for each qubit count from 3 to 26, selecting the densest QFT variant that fits the IQS index budget of 603 and disabling final QFT swaps. We run only split SQL execution on PostgreSQL, DuckDB, and SQLite, sweeping Docker memory caps of 16, 8, and 4~GB. Each engine-cap-circuit tuple is executed with one warm-up run followed by five timed runs, and the runtime tables report the median of the five timed runs.

% All containers use a one-CPU quota, swap is disabled by setting Docker
% \texttt{--memory-swap} equal to the memory cap, and host page cache is not
% explicitly dropped between runs. DuckDB uses one thread and sets
% \texttt{memory\_limit} to the limit minus a 1024~MB headroom pad (3072, 7168,
% and 15360~MB for the 4, 8, and 16~GB caps, respectively). SQLite uses a
% 64~MB cache budget for both main and temporary schemas. PostgreSQL uses the
% tuned PostgreSQL~12.22 container, with the server container capped at the
% experiment memory limit and \texttt{temp\_file\_limit} set to 64~GB. Spill is
% measured as \texttt{postgres\_explain\_temp\_written} for PostgreSQL and as
% \texttt{proc\_io\_write\_bytes} for DuckDB and SQLite.
% For runtime, we report the median runtime of five repeated runs.

% To get the spill results, for PostgreSQL we use \texttt{postgres\_explain\_temp\_written}, and \texttt{proc\_io\_write\_bytes} for DuckDB and SQLite.
% \section{Results on Out-of-Core Evaluation}
% % \section{Memory-limit sensitivity for QFT at 26 qubits}
% \label{app:qft-cap-sensitivity}

% We summarize the memory-limit sensitivity of split-query execution for dense QFT circuits from 3 to 26 qubits on PostgreSQL, DuckDB, and SQLite, \rh{here add which container you used} using one container CPU and memory caps of 4, 8, and 16~GB. Each configuration is run once as a warmup and five times for measurement; we report medians over the five timed runs. The dataset contains 1{,}291 rows, compared with an expected full-success grid of 1{,}296; 1{,}290 rows complete successfully and one is terminated by an OOM kill. The only failure is DuckDB at 26 qubits under the 4~GB cap. In that row, the recorded gate count is truncated at 51, indicating that execution stopped before the full circuit completed. \rh{add ``refer to Table.11-13'' at the right place.}

% \para{Postgres.}
% Postgres exhibits a consistently heavy-spill regime across all three caps. As qubit count increases from 23 to 26, spill grows monotonically from 2.32~GB to 23.49~GB, and this progression is effectively unchanged across the 4, 8, and 16~GB settings. Runtime improves at higher caps for the largest instances, but the engine continues to rely on substantial externalization even at 16~GB. For dense QFT contractions, Postgres therefore behaves as a stable out-of-core engine rather than transitioning into a light-spill regime.

% \para{DuckDB.}
% DuckDB shows the sharpest change between heavy- and light-spill behavior. Up to 24 qubits, spill remains negligible across all caps. At 25 qubits, the 4~GB setting jumps to 7.68~GB of spill while the 8~GB and 16~GB settings remain near zero. At 26 qubits, the 4~GB run fails, the 8~GB run completes with heavy spill (9.69~GB), and the 16~GB run completes with near-zero spill. Thus, as qubit count grows, DuckDB moves from a light-spill regime into a threshold region where successful execution depends strongly on the available memory cap.

% \para{SQLite.}
% SQLite shows a gradual but persistent spill pattern. From 23 to 26 qubits, spill rises steadily from about 0.51~GB to about 6.1~GB across all three caps, while runtime also increases monotonically. Unlike DuckDB, SQLite does not exhibit a sharp transition between low- and high-spill behavior; instead, it scales through the dense QFT workload by relying on moderate out-of-core execution throughout. This makes SQLite slower than Postgres at the largest sizes, but more robust than DuckDB at 4~GB.

% \para{Results and relevance to \sys.}
% Taken together, these results show that dense QFT contractions expose three distinct engine-level spill profiles: stable heavy spill for Postgres, threshold-sensitive spill for DuckDB, and gradual moderate spill for SQLite. The qubit-by-qubit evolution in Tables~\ref{tab:qft-26q-runtime-evolution} and~\ref{tab:qft-26q-spill-evolution} makes these differences explicit. More broadly, this experiment illustrates a central contribution of \sys: it enables systematic study of out-of-core execution by exposing workload- and engine-dependent spill behavior alongside runtime and circuit features.

% \begin{table}[t]
%   \centering
%   \caption{26-qubit QFT results under split execution for three memory caps. Runtimes are medians over five timed runs after one warmup. Spill is measured by \texttt{postgres\_explain\_temp\_written} for Postgres and by \texttt{proc\_io\_write\_bytes} as a spill proxy for DuckDB and SQLite. DuckDB at 4~GB is OOM-killed before finishing the full circuit; the recorded \texttt{num\_gates}=51 is therefore truncated and is not a completed runtime observation.}
%   \label{tab:qft-26q-cap-summary}
%   \footnotesize
%   \begin{tabularx}{\columnwidth}{@{}lXXX@{}}
%     \toprule
%     \textbf{Cap} & \textbf{Postgres} & \textbf{DuckDB} & \textbf{SQLite} \\
%     \midrule
%     4~GB & success; 223.5~s; 23.49~GB spill & OOM kill; truncated at 51 gates & success; 423.2~s; 6.14~GB spill \\
%     8~GB & success; 214.5~s; 23.49~GB spill & success; 199.9~s; 9.69~GB spill & success; 392.3~s; 5.91~GB spill \\
%     16~GB & success; 154.6~s; 23.49~GB spill & success; 189.5~s; $\sim$0~GB spill & success; 405.5~s; 6.07~GB spill \\
%     \bottomrule
%   \end{tabularx}
% \end{table}

% \begin{table}[t]
%   \centering
%   \caption{Runtime evolution (seconds) for the high-qubit QFT region under split execution. Each entry is the median over five timed runs.}
%   \label{tab:qft-26q-runtime-evolution}
%   \footnotesize
%   \setlength{\tabcolsep}{4pt}
%   \begin{tabularx}{\columnwidth}{@{}llXXXX@{}}
%     \toprule
%     \textbf{Engine} & \textbf{Cap} & \textbf{23q} & \textbf{24q} & \textbf{25q} & \textbf{26q} \\
%     \midrule
%     Postgres & 4~GB  & 21.9 & 46.5 & 101.6 & 223.5 \\
%     Postgres & 8~GB  & 17.8 & 46.6 & 74.8  & 214.5 \\
%     Postgres & 16~GB & 22.2 & 46.5 & 74.8  & 154.6 \\
%     DuckDB   & 4~GB  & 21.2 & 45.1 & 104.5 & --    \\
%     DuckDB   & 8~GB  & 21.2 & 43.7 & 91.4  & 199.9 \\
%     DuckDB   & 16~GB & 21.5 & 43.5 & 91.7  & 189.5 \\
%     SQLite   & 4~GB  & 43.3 & 92.8 & 188.9 & 423.2 \\
%     SQLite   & 8~GB  & 46.3 & 96.0 & 180.1 & 392.3 \\
%     SQLite   & 16~GB & 42.1 & 92.5 & 186.8 & 405.5 \\
%     \bottomrule
%   \end{tabularx}
% \end{table}

% \begin{table}[t]
%   \centering
%   \caption{Spill evolution (GB) for the high-qubit QFT region under split execution. Postgres uses \texttt{postgres\_explain\_temp\_written}; DuckDB and SQLite use \texttt{proc\_io\_write\_bytes} as a spill proxy.}
%   \label{tab:qft-26q-spill-evolution}
%   \footnotesize
%   \setlength{\tabcolsep}{4pt}
%   \begin{tabularx}{\columnwidth}{@{}llXXXX@{}}
%     \toprule
%     \textbf{Engine} & \textbf{Cap} & \textbf{23q} & \textbf{24q} & \textbf{25q} & \textbf{26q} \\
%     \midrule
%     Postgres & 4~GB  & 2.32 & 4.64 & 9.82 & 23.49 \\
%     Postgres & 8~GB  & 2.32 & 4.64 & 9.82 & 23.49 \\
%     Postgres & 16~GB & 2.32 & 4.64 & 9.82 & 23.49 \\
%     DuckDB   & 4~GB  & 0.00 & 0.00 & 7.68 & --    \\
%     DuckDB   & 8~GB  & 0.00 & 0.00 & 0.00 & 9.69  \\
%     DuckDB   & 16~GB & 0.00 & 0.00 & 0.00 & 0.00  \\
%     SQLite   & 4~GB  & 0.51 & 1.22 & 2.71 & 6.14  \\
%     SQLite   & 8~GB  & 0.51 & 1.24 & 2.62 & 5.91  \\
%     SQLite   & 16~GB & 0.51 & 1.22 & 2.63 & 6.07  \\
%     \bottomrule
%   \end{tabularx}
% \end{table}

% \begin{figure}[t]
%   \centering
%   \IfFileExists{images/qft_26q_caps_runtime_spill.pdf}{%
%     \includegraphics[width=\columnwidth]{images/qft_26q_caps_runtime_spill.pdf}%
%   }{%
%     \includegraphics[width=\columnwidth]{images/qft_26q_caps_runtime_spill.png}%
%   }
%   \caption{Memory-limit sensitivity of dense QFT contractions at 26 qubits under split execution. Bars are grouped by engine and colored by memory limit (4, 8, and 16~GB). The left panel reports median runtime over five timed runs after one warmup, and the right panel reports spill volume. The disk marker above each successful bar summarizes normalized spill volume relative to the largest successful spill observation in this experiment. For Postgres, spill is measured by \texttt{postgres\_explain\_temp\_written}; for DuckDB and SQLite, spill is approximated by \texttt{proc\_io\_write\_bytes}. DuckDB at 4~GB is OOM-killed before completing the circuit and records a truncated gate count (\texttt{num\_gates}=51); it is therefore shown as a failure marker rather than a completed bar.}
%   \label{fig:qft-26q-cap-sensitivity}
% \end{figure}



% \medskip
% \para{Runtime performance}
% Table~\ref{tab:qft-26q-runtime-evolution} reports the runtime of SQL queries
% generated from larger, high-qubit-count circuits across RDBMSs and memory limits,
% where spilling may occur under memory limit. PostgreSQL and DuckDB have
% comparable runtimes up to 24 qubits. At 25 and 26 qubits, memory pressure becomes
% dominant: PostgreSQL benefits from a larger memory limit, whereas DuckDB exhibits
% a sharp transition from successful execution to out-of-memory failure at 26 qubits
% under the 4~GB limit. SQLite is consistently slower, but its runtime increases
% smoothly and it completes under all memory limits.

% \begin{table}[t]
%   \centering
%   \caption{Runtime evolution in seconds. A dash indicates the OOM configuration.}
%   \label{tab:qft-26q-runtime-evolution}
%   \footnotesize
%   \setlength{\tabcolsep}{4pt}
%   \begin{tabular}{@{}llccc@{}}
%     \toprule
%     \textbf{Engine} & \textbf{Qubits} & \textbf{4~GB} & \textbf{8~GB} & \textbf{16~GB} \\
%     \midrule
%     \multirow{4}{*}{PostgreSQL}
%       & 23 & 21.9  & 17.8  & 22.2  \\
%       & 24 & 46.5  & 46.6  & 46.5  \\
%       & 25 & 101.6 & 74.8  & 74.8  \\
%       & 26 & 223.5 & 214.5 & 154.6 \\
%     \midrule
%     \multirow{4}{*}{DuckDB}
%       & 23 & 21.2  & 21.2  & 21.5  \\
%       & 24 & 45.1  & 43.7  & 43.5  \\
%       & 25 & 104.5 & 91.4  & 91.7  \\
%       & 26 & --    & 199.9 & 189.5 \\
%     \midrule
%     \multirow{4}{*}{SQLite}
%       & 23 & 43.3  & 46.3  & 42.1  \\
%       & 24 & 92.8  & 96.0  & 92.5  \\
%       & 25 & 188.9 & 180.1 & 186.8 \\
%       & 26 & 423.2 & 392.3 & 405.5 \\
%     \bottomrule
%   \end{tabular}
% \end{table}

% \medskip
% \para{Spill performance}
% Table~\ref{tab:qft-26q-spill-evolution} reports spill volume.
% % \footnote{We measure PostgreSQL spill as the temporary blocks written reported by
% % \texttt{EXPLAIN (ANALYZE, BUFFERS)}, denoted
% % \texttt{postgres\_explain\_temp\_written}; for DuckDB and SQLite, we approximate
% % spill using the per-query delta of Linux
% % \texttt{/proc/\textless pid\textgreater/io}'s \texttt{write\_bytes}, denoted
% % \texttt{proc\_io\_write\_bytes}. See
% % \url{https://www.postgresql.org/docs/current/sql-explain.html} and
% % \url{https://docs.kernel.org/filesystems/proc.html#proc-pid-io-display-the-io-accounting-fields}.}
% PostgreSQL spill grows
% monotonically with qubit count and is essentially unchanged across memory limits,
% suggesting a stable heavy-spill execution mode. DuckDB shows a different behavior:
% it has almost no spill until a threshold is crossed, after which execution either spills
% heavily or fails. SQLite spills gradually as qubit count increases and remains stable
% across memory limits. These differences are precisely the kind of engine-specific
% behavior that \sys makes visible for out-of-core quantum simulation workloads.

% \medskip
% Moreover, now we focus on the largest QFT circuit (26 qubits), PostgreSQL completes all three memory limits, but remains spill-heavy: it writes 23.49~GB of temporary data even at 16~GB. DuckDB is highly memory-threshold sensitive: it fails at 4~GB, completes at 8~GB with 9.69~GB of spill, and completes at 16~GB with near-zero spill. SQLite completes all three memory limits with moderate spill, but has the highest runtime at 26 qubits. Thus, the engines expose three distinct out-of-core
% profiles: robust heavy spilling (PostgreSQL), threshold-sensitive spilling (DuckDB),
% and robust moderate spilling (SQLite). These differences suggest that for future work on further improving
% RDBMS-backed simulation for large quantum circuits, we can investigate
% engine-specific
% query tuning, physical design, and spill-aware execution strategies.

% Through these results, we observe that for large circuits that exceed memory, RDBMS engines can
% still complete simulation through out-of-core execution. This
% highlights a distinct advantage of database-backed simulation and
% enables follow-up work on spill-aware SQL plans and memory-
% sensitive query optimization.

% \begin{table}[t]
%   \centering
%   \caption{Spill evolution in GB. A dash indicates the OOM.}
%   \label{tab:qft-26q-spill-evolution}
%   \footnotesize
%   \setlength{\tabcolsep}{4pt}
%   \begin{tabular}{@{}llccc@{}}
%     \toprule
%     \textbf{Engine} & \textbf{Qubits} & \textbf{4~GB} & \textbf{8~GB} & \textbf{16~GB} \\
%     \midrule
%     \multirow{4}{*}{PostgreSQL}
%       & 23 & 2.32  & 2.32  & 2.32  \\
%       & 24 & 4.64  & 4.64  & 4.64  \\
%       & 25 & 9.82  & 9.82  & 9.82  \\
%       & 26 & 23.49 & 23.49 & 23.49 \\
%     \midrule
%     \multirow{4}{*}{DuckDB}
%       & 23 & 0.00 & 0.00 & 0.00 \\
%       & 24 & 0.00 & 0.00 & 0.00 \\
%       & 25 & 7.68 & 0.00 & 0.00 \\
%       & 26 & --   & 9.69 & 0.00 \\
%     \midrule
%     \multirow{4}{*}{SQLite}
%       & 23 & 0.51 & 0.51 & 0.51 \\
%       & 24 & 1.22 & 1.24 & 1.22 \\
%       & 25 & 2.71 & 2.62 & 2.63 \\
%       & 26 & 6.14 & 5.91 & 6.07 \\
%     \bottomrule
%   \end{tabular}
% \end{table}

 


% \section{Out-of-Core Evaluation}
% \label{app:qft-cap-sensitivity}

% In this section, we explore \emph{out-of-core}
% execution for DBMS-backed SQL simulation, and try to answer the following research questions: 
% Q3: Can RDBMS-based simulation complete workloads that no
% longer fit in memory?




% We have  generate a set of large circuits, one circuit for each qubit count from one to 45. 
% simulated  in Section~\ref{app:ooc-expander}. With the memory limit, Qiskit Aer cannot finish the simulation. In contract, RDBMSs spill to secondary storage to complete execution under memory limit, demonstrating a clear case \emph{when} it is benifitial to use RDBMS as a simulation backend. 

% Through this experiment, we observed that the effective
%  simulation
% of  RDBMS is becuase they exploit the sparsity of the input circuits. Therefore, we designed a second experiment in Sec.~\ref{app:ooc-inferq-sparsity}, in which we vary the sparisty of the generated circuits via \sys, and observe  sparsity, intermediate resutls, and spill behavior.
 


% \subsection{Large Circuits Simulation}
% \label{app:ooc-expander}
% \begin{figure}[t]
%     \centering
%     \includegraphics[width=1.0\linewidth]{images/expanderfig3_spill_vs_cte.pdf}
%     % all gb results stores in allgb_ prefix image, no explanation so omitted. 
%     \caption{Out-of-core behavior (Median Spill for 5 trials of query run) on 40--45 qubit Expander circuits under 16~GB caps. Filled points (tot : total) denote total intermediate relation size across the CTE chain; open points (lrg : largest) denote the maximum single-step intermediate size.}
%     \label{fig:expander_ooc}
% \end{figure}

% \para{Experiment setup.}
% we generate a set of circuits with the same structure and sparsity but varying number of qubits $N$, each circuit for each qubit count from one to 45. This set of circuits are relatively sparse. 
% we generate a set of circuits with the same structure and sparsity but varying number of qubits $N$, each circuit for each qubit count from one to 45. This set of circuits are relatively sparse. 
% This set of sparse circuits are presentative  in practice: stabilizer, amplitude amplification, and and sparse Hamiltonians and sparse linear systems.
% Moreover,
% even when the final state is dense, intermediate tensors and relations during contraction
% can be much smaller than $2^N$.  
% %
% We simulate a set of circuits with the same structure but varying number of qubits $N$. The bigger value $N$ is, the larger circuit is and running the SQL query for simulating this circuit would require more memory.  Quantum-specific details about these circuits are provided in Section~\ref{app:expander-details}. 
% %
% We run PostgreSQL, DuckDB, and SQLite under  
% memory caps of 16~GB and record (i) median spill volume over 5 runs per circuit query and (ii) the maximum and total
% intermediate relation sizes across the CTE chain during SQL execution. 
% The detailed configuration for each database engine for the out-of-core experiments is in Section~\ref{app:spill-instrumentation}.


% \para{Results.} Under 16~GB limit, Qikit Aer failed to complete the simulation from circuits with 30 qubits. In contrast, all RDBMS have finished all simulations.
% Figure~\ref{fig:expander_ooc} summarizes spill behavior and intermediate sizes for the
% 40--45 qubit circuits. 

% First, why did Qikit Aer fail? Recall in Sec. 1 and 2.1 we have explained that for $N$-qubit state it leads to $2^n$ complex amplitudes. If we store each of these complex numbers in double
% precision, even if we ignore overhead, 
% this is $16 \cdot 2^N$ bytes, which reaches roughly 16~GB at
% $n{=}30$ and grows exponentially thereafter. 
% This is how Qiskit Aer an exact simulation, known as the \texttt{statevector} method. 
% So, for this set of circuits with 40 to 45 qubits,  where a Qiskit Aer
% would require approximately 16--512~TB. 
% % requires storing $2^n$ complex amplitudes.
% This is why in this experiments Qikit Aer fail. 
% % That is, 
% % It  cannot execute circuits once $2^n$ exceeds available RAM.
% In contrast, RDBMSs stores amplitudes as tuples and can exploit the sparsity: if only a small subset of basis states have non-zero amplitudes,
% the representation size scales with the number of non-zero amplitudes rather than grows exponentially with qubit count ($2^N$).


% Now we anlyze Figure~\ref{fig:expander_ooc}. Across engines, intermediate relation sizes remain in the megabyte (MB) range, consistent with the RDBMSs'  sparsty support of the state representation.\footnote{We also observe that different spilling behvoirs among engines.  DuckDB
% consistently spills the least, while SQLite and PostgreSQL are comparable on several
% circuits. This is most likely because the different inner mechnism of these engines. However, we should also mention that for PostgreSQL we cancompute exact temporary spill as this is directly supported. From the public availble documentation, DuckDB and SQLite do not provide such support so we can only  use a conservative operating system level proxy for them. We have added all details about this in Section~\ref{app:spill-instrumentation}.}
% We see that Spill volume is often orders of magnitude larger than the final and intermediate
% relation sizes. This might be becuase additional overhead required by join and aggregate operators
% (e.g., hash tables, sorting or grouping state, and temporary structures) and repeated
% materialization of intermediates. 
% Now let's look at the Spearman correlation values of  maximum single-step intermediate size (denoted as $r(lrg)$) and  total intermediate relation
% size across the CTE chain (denoted as $r(tot)$) in Figure~\ref{fig:expander_ooc}. 
% Spill correlates more strongly
% with the total intermediate size across all CTEs than with the single largest CTE
% (Spearman correlation higher for total-CTE than max-CTE in all settings). This indicates that for simulation over this set of sparse circuits, spilling volume is more correlated with all CTEs than with the single largest CTE. For future work, it might worth investating more sophsicated optimiazation, e.g., materialize and cache CTEs that are repeatedly used. 
% % Varying the memory limitcauses only modest changes in intermediate relation
% % sizes, indicating that the DBMS execution primarily benefits from sparse  
% % representation, while out-of-core externalization absorbs operator workspace peaks.


% \para{Takeaways.}
% This experiment demonstrates when DBMS-based simulation is uniquely useful for quantum circuit simulation:
% exact simulation for large-$N$ circuits that exceed the in-memory Qiskit Aer engine's capbility (dense
% statevector representation), by combining sparse tuple-level representation with out-of-core
% execution. At the same time, spill volume is correlated with relational operators and intermediate
% relation growth, which motivates query optimizations (join ordering, operator choice,
% and materialization control) to reduce externalization overhead.
\section{Out-of-Core Evaluation}
\label{app:qft-cap-sensitivity}

In this section, we study \emph{out-of-core} execution for DBMS-backed SQL simulation and
address the following question:

\smallskip
\noindent\textbf{Q3:} Can an RDBMS backend complete simulation queries once the workload no
longer fits in main memory?

\smallskip
To answer Q3, we conduct two experiments. First, in
Section~\ref{app:ooc-expander}, we generate a family of large circuits with an increasing number of qubits up to 45. Under a fixed memory limit, Qiskit Aer cannot complete the simulation, whereas RDBMS engines can spill
temporary data to secondary storage and complete execution, demonstrating a clear case
\emph{when} an RDBMS backend is beneficial. Second, in
Section~\ref{app:ooc-inferq-sparsity}, we use \sys to generate circuits with controlled
variation in sparsity and study how sparsity, intermediate results, and spill behavior
interact.

\subsection{Large Circuits Simulation}
\label{app:ooc-expander}



\para{Experiment setup}
We generate a family of circuits with a fixed structure and similar sparsity, while
increasing the qubit count $N$ (one circuit for each $N$ from 1 to 45). These circuits
are designed so that the simulation state can be represented compactly as a relation of
non-zero entries. These circuits are representative in practice and are a common scenario for structured and sparse quantum circuit simulation workloads (e.g., stabilizer, amplitude amplification, and sparse linear systems). 
Quantum-specific construction details are provided in Section~\ref{app:expander-details}.

We execute the generated SQL queries on PostgreSQL, DuckDB, and SQLite under a 16~GB memory cap. In SQL-based simulation, the tensor-contraction sequence is written as a chain of Common Table Expressions (CTEs), each encoding one contraction step (join + aggregate). 
For each query, we record (i) median spill volume over five runs and (ii) the
maximum and total intermediate relation sizes across the CTE chain. The engine-specific
spill configuration is described in
Section~\ref{app:spill-instrumentation}.


\para{Results}
Under the 16~GB limit, Qiskit Aer fails to complete exact simulation starting at 30 qubits.
This is expected. Recall in Sec. 1 and 2.1 we have explained that 
an $N$-qubit state requires $2^N$
complex entries. If we store each complex entry in double precision, this is approximately $16 \cdot 2^N$ bytes,
which reaches roughly 16~GB at $N{=}30$ (ignoring runtime overhead) and grows
exponentially thereafter. For the 40--45 qubit circuits evaluated here, Qiskit would require approximately 16--512~TB (this is known as
statevector in Qiskit for exact simulation). 

In contrast, all three RDBMS engines
complete all runs because they store the state as tuples and can exploit sparsity: storage
scales with the number of non-zero entries rather than with $2^N$.



Figure~\ref{fig:expander_ooc} summarizes out-of-core behavior for the 40--45 qubit circuit simulation under the 16~GB limit.
Across engines, intermediate relation sizes remain in the megabyte range, consistent with RDBMSs' support for storing sparse tensors.\footnote{We observe different spill behavior across DBMS
engines. DuckDB spills the least in this experiment, while SQLite and PostgreSQL are
comparable on several circuits. Note that PostgreSQL exposes query-level temporary I/O
counters, so we can compute spill volume exactly. DuckDB and SQLite do not expose a
stable, query-level temp-bytes-written counter in the same way, so we use a conservative
operating-system level proxy for them. Details are provided in
Section~\ref{app:spill-instrumentation}.}
At the same time, total spill volume is often orders of magnitude larger than the sizes of the
intermediate relations. This is likely due to the overhead during joins and
group-bys (e.g., hash tables, sorting or grouping state, and temporary structures), as
well as repeated materialization of intermediates along the CTE chain.




We further inspect the Spearman correlation values between spill volume and the maximum
single-step intermediate size ($r(\texttt{lrg})$) and the total intermediate size across
all CTEs ($r(\texttt{tot})$). In all three engines, $r(\texttt{tot}) > r(\texttt{lrg})$,
indicating that spill is more strongly associated with cumulative intermediate volume than
with the peak single-step intermediate. This suggests that optimization opportunities
include reducing repeated intermediate materialization and improving reuse, for example by
caching selected intermediates when they are used multiple times.

\begin{figure}[t]
    \centering
    \includegraphics[width=1.0\linewidth]{images/expanderfig3_spill_vs_cte.pdf}
    \caption{Spill volume of RDBMS engines (PostgreSQL, DuckDB, SQLite). Filled points (\texttt{tot}) denote total intermediate relation size across the CTE chain; open points (\texttt{lrg}) denote the maximum single-step intermediate size.}
    \label{fig:expander_ooc}
\end{figure}
\begin{figure}[t]
    \centering
    \includegraphics[width=1.0\linewidth]{images/fig4a_walltime_vs_qubits.pdf}
    \caption{Runtime of RDBMS engines (PostgreSQL, DuckDB, SQLite).}
    \label{fig:expander_walltime_qubits}
\end{figure}

\begin{figure}[t]
    \centering
    \includegraphics[width=1.0\linewidth]{images/fig4b_walltime_vs_spill.pdf}
    \caption{Runtime versus spill volume for RDBMS engines (PostgreSQL, DuckDB, SQLite).}
    \label{fig:expander_walltime_spill}
\end{figure}


\begin{figure*}[t]
    \centering
    \includegraphics[width=0.55\linewidth]{images/fig2_spill_vs_qubits.pdf}
    \caption{Spill volume with varying numbers of qubits (log scale).}
    \label{fig:dbspillqubits_repeat}
\end{figure*}
\begin{figure*}[t]
    \centering
    \includegraphics[width=0.6\linewidth]{images/fig7_cte_vs_density.pdf}
    \caption{Maximum and total intermediate CTE sizes against output state sparsity (log scale).}
    \label{fig:dbspillctesparse_repeat}
\end{figure*}

Figure~\ref{fig:expander_walltime_qubits} reports query runtime. 
Runtime increases rapidly with the number of qubits for all engines. PostgreSQL is consistently the fastest in this experiment, DuckDB is second,
and SQLite is the slowest. Figure~\ref{fig:expander_walltime_spill} shows runtime and
spill volume are related, indicating that externalization overhead is a major
reason for end-to-end query time once operator workspaces spill to disk.

\medskip
  % \vspace{0.2cm}
\noindent
\setlength{\fboxsep}{6pt}
\fcolorbox{black!15}{black!4}{%
  \parbox{\dimexpr\linewidth-2\fboxsep-2\fboxrule\relax}{
\para{Takeaways}
This experiment demonstrates \emph{when} DBMS-backed simulation is particularly useful:
it enables exact simulation for large circuits that exceed the in-memory limits of
Qiskit Aer by combining a sparse tuple-level representation with out-of-core
execution. At the same time, spill volume and runtime are closely tied to relational
operator workspace and intermediate relation growth, which motivates RDBMS
optimization opportunities such as join ordering, operator selection, and materialization
control to reduce externalization overhead.
}}

\begin{figure*}[t]
    \centering
    \includegraphics[width=0.5\linewidth]{images/fig8_walltime_vs_spill.pdf}
     \vspace{-0.3cm}
    \caption{Execution time versus spill volume for sampled InferQ circuits (log-log).}
    \label{fig:dbwalltimespill_repeat}
    \vspace{-0.4cm}
\end{figure*}
\subsection{Experiment on Varying Sparsity}
\label{app:ooc-inferq-sparsity}

\begin{table*}[!ht]
  \centering
  \caption{Ablation study for \emph{time-optimal method} prediction
    (80\% - 20\% train-test split).
    \textbf{Bold}: best overall result (XGB, full feature set);
    \underline{underline}: second best (RF, full feature set).
    Acc: accuracy; F1: macro F1-score.}
  \label{tab:ablation_time}
  \scriptsize
  \setlength{\tabcolsep}{3pt}
  \renewcommand{\arraystretch}{1.05}
  \begin{tabular}{l cc cc cc cc cc}
    \toprule
    & \multicolumn{2}{c}{\textbf{LR}}
    & \multicolumn{2}{c}{\textbf{SVM}}
    & \multicolumn{2}{c}{\textbf{DT}}
    & \multicolumn{2}{c}{\textbf{RF}}
    & \multicolumn{2}{c}{\textbf{XGB}} \\
    \cmidrule(lr){2-3}\cmidrule(lr){4-5}\cmidrule(lr){6-7}
    \cmidrule(lr){8-9}\cmidrule(lr){10-11}
    \textbf{Feature Set}
      & Acc & F1
      & Acc & F1
      & Acc & F1
      & Acc & F1
      & Acc & F1 \\
    \midrule
    SQL
      & 0.79 & 0.54
      & 0.80 & 0.57
      & 0.91 & 0.70
      & 0.93 & 0.75
      & 0.92 & 0.75 \\
    Static
      & 0.80 & 0.55
      & 0.80 & 0.55
      & 0.90 & 0.68
      & 0.92 & 0.72
      & 0.92 & 0.74 \\
    Static + Graph
      & 0.84 & 0.60
      & 0.81 & 0.56
      & 0.90 & 0.66
      & 0.92 & 0.73
      & 0.92 & 0.74 \\
    Static + SQL
      & 0.86 & 0.64
      & 0.86 & 0.65
      & 0.93 & 0.76
      & 0.94 & 0.78
      & 0.94 & 0.80 \\
    Static + Dynamic
      & 0.86 & 0.66
      & 0.86 & 0.65
      & 0.92 & 0.73
      & 0.94 & 0.77
      & 0.94 & 0.79 \\
    Static + Graph + SQL
      & 0.87 & 0.67
      & 0.84 & 0.62
      & 0.92 & 0.74
      & 0.94 & 0.78
      & 0.93 & 0.78 \\
    Static + Graph + SQL + Dyn.
      & 0.89 & 0.70
      & 0.88 & 0.69
      & 0.94 & 0.80
      & \underline{0.95} & \underline{0.83}
      & \textbf{0.95} & \textbf{0.84} \\
    \bottomrule
  \end{tabular}
\end{table*}

The previous experiment indicates that sparse representations can make DBMS-backed
simulation more useful under a memory limit.   
We therefore conduct a second experiment on
InferQ-generated workloads to study how spill volume, intermediate CTE sizes, and query
runtime relate to the sparsity of a circuit.
This is possible because InferQ can generate circuits
with controlled variation in sparsity.

We first explain the sparsity of the final output state of a circuit simulation.
For an $n$-qubit circuit, the output is a length-$2^n$ amplitude vector
$\mathbf{p}$. We define \emph{output density} as the fraction of non-zero entries in the vector $\mathbf{p}$ against the vector size $|\mathbf{p}|$:
\[
\mathrm{density}(\mathbf{p}) \;=\; \frac{\left|\{\, i \mid p_i \neq 0 \,\}\right|}{|\mathbf{p}|}.
\]
A smaller density indicates fewer non-zero tuples in the DBMS representation, while a
density close to 1 indicates a near-dense output.

\medskip
\para{Experiment setup}
We construct a sample of 17 InferQ circuits\footnote{\url{https://github.com/InfiniData-Lab/InferQ/blob/main/analysis/sample_csvs/final_sparse_sample.csv}}
spanning three density bands: Sparse $[0,0.33)$, Medium $[0.33,0.67)$, and Dense
$[0.67,1.0]$. The circuits span a number of qubits from 20 to 23. 
We execute the
SQL simulation on PostgreSQL, DuckDB, and SQLite under a 16~GB memory cap. For each run,
we record spill volume, wall-clock runtime, and the maximum and total intermediate CTE
sizes observed over the CTE chain.

\medskip
\para{Results}
Figure~\ref{fig:dbspillqubits_repeat} plots spill volume against the number of qubits. Across
all engines, circuits in the Sparse band tend to incur the least spill, and PostgreSQL
often shows negligible spill for the sparsest cases. This matches the intuition that a
sparse output can be represented as a smaller relation, reducing memory pressure and
temporary-workspace needs.

From Figure~\ref{fig:dbspillqubits_repeat}, we observe that output density alone does not determine spill volume. For PostgreSQL and SQLite, several
Medium-density circuits spill more than some Dense circuits at similar qubit counts. This indicates that spill is strongly influenced by
\emph{intermediate} results produced during query execution, which depend on plan choices
such as join ordering and materialization behavior, rather than being determined solely
by the final output density.




Figure~\ref{fig:dbspillctesparse_repeat} reports the maximum and total intermediate CTE sizes as
a function of output density. The correlations are consistently
positive across engines, and the maximum-intermediate correlation is slightly stronger
than the total-intermediate correlation. At the same time, intermediate sizes vary across
DBMS engines for the same circuit, which is expected because optimizers may choose
different join orders and physical operators, leading to different intermediate growth.

Finally, Figure~\ref{fig:dbwalltimespill_repeat} shows runtime versus spill volume. PostgreSQL and
SQLite show a very strong positive association between execution time and spill (rank
correlations shown in the figure), and DuckDB exhibits a weaker but
still positive association. This result indicates that externalization cost is a dominant
reason for end-to-end runtime once operators spill to disk. 

\medskip
  \vspace{0.2cm}
\noindent
\setlength{\fboxsep}{6pt}
\fcolorbox{black!15}{black!4}{%
  \parbox{\dimexpr\linewidth-2\fboxsep-2\fboxrule\relax}{
\para{Takeaways}
It shows that DBMS out-of-core behavior is governed primarily by
intermediate relation growth and operator workspace, which depend on workload structure
and DBMS plan choices, rather than only on the density of the final output. It motivates
optimizing SQL-based simulation using standard DBMS techniques such as improving join
ordering, controlling materialization, and reducing intermediate blow-up. It also
suggests that spill-aware cost models are useful for predicting when a DBMS   will
benefit from out-of-core execution and when externalization overhead will dominate in future work.
}}


% \subsection{Experiment on Varying Sparsity Circuits}
% \label{app:ooc-inferq-sparsity}

% % \para{Why this experiment.}
% The previous
% Experiment help us relaize sparisty matters for RDBMS's out-of-core perofrmance.  We next conduct a second experiment, in which we have circuits with varying sparisty, more specificlly, how spill volume and intermediate CTE sizes
% relate to output-state sparsity. This is possible because InferQ can generate circuits
% with controlled variation in sparsity.

% \para{Experiment setup}
% We construct a sample of 17 InferQ circuits \footnote{\url{https://anonymous.4open.science/r/InferQ-ADEF/analysis/sample_csvs/final_sparse_sample.csv}} spanning three output-sparsity bands. Sparsity
% is defined as the fraction of non-zero amplitudes in the statevector $\mathbf{p}$:
% \[
% \mathrm{sparsity}(\mathbf{p}) \;=\; \frac{\left|\{\, i \mid p_i \neq 0 \,\}\right|}{|\mathbf{p}|}.
% \]
% We use three bands: Sparse $[0,0.33)$, Medium $[0.33,0.67)$, and Dense $[0.67,1.0]$. The
% circuits span \texttt{num\_qubit} from 20 to 23. While this range is within the
% capabilities of in-memory statevector simulators, it is sufficient to observe meaningful
% differences in spill and intermediate sizes across DBMS engines. We execute the SQL
% simulation on PostgreSQL, DuckDB, and SQLite under a 16~GB memory limitand record spill
% volume as well as the maximum and total intermediate CTE sizes during query execution.


% \para{Results.}
% Figure~\ref{fig:dbspillqubits_repeat} shows the results.  
% % DuckDB spills the least at these qubit counts,
% % while PostgreSQL and SQLite spill one to two orders of magnitude more than DuckDB in several cases.
% % This is consistent with the larger-qubit QFT study, where DuckDB begins to spill more
% % heavily beyond 25 qubits (Table~\ref{tab:qft-26q-spill-evolution}). 
% Across all engines,
% for circuits in the Sparse band, all engines tend to have the lowest spill, with PostgreSQL showing minimal and often
% negligible spill on the most sparse circuits.

% We can see that the sparsity of the final output state alone does not determine spill volume. For PostgreSQL and SQLite, several
% Medium-sparsity circuits spill more than some Dense circuits at similar qubit counts
% (Figure~\ref{fig:dbspillqubits_repeat}), which indicates that spill is influenced by intermediate
% states and the resulting intermediate relations produced by the execution plan, not only
% by the sparsity of the final output state.

% Figure~\ref{fig:dbspillctesparse_repeat} reports maximum and total intermediate CTE sizes versus
% output sparsity. Both metrics vary substantially across circuits, and in many cases the
% Medium band produces intermediate CTE sizes comparable to or larger than the Dense band.
% Since each DBMS may choose different join orders and execution strategies, intermediate
% CTE sizes also differ across engines for the same circuit. Overall, sparsity is
% monotonically correlated with intermediate CTE sizes (Spearman coefficient above 0.7),
% with the correlation slightly stronger for maximum CTE size than for total CTE size.

% \medskip
% \para{Takeaways.}
% This experiment help us understand better the relationship between RDBMSs' out-of-core performance, and intermediate
% relation growth during contraction, which depends on both circuit structure and DBMS plan
% choices (e.g., join ordering and materialization behavior), rather than only the sparsity
% of the final state. It motivates future work on reducing intermediate blow-up through
% better contraction ordering and plan selection. It might also be possible to try out circuit-level transformations
% that compress intermediate states to better leverage DBMS backends.
\section{Ablation Study of InferQ Feature groups}
\label{appendix:abl}


% %  Table: Memory-optimal method prediction (best_mem_method)
% %  Shows best model per feature set. Bold = best in column.

% \begin{table}[t]
%   \centering
%   \caption{%
%     Ablation study for \emph{memory-optimal method} prediction
%     (train=6{,}164 / test=1{,}541; class split 58.4\% / 41.6\%).
%     Each row reports the best-performing model for that feature set.
%   }
%   \label{tab:ablation_mem}
%   \small
%   \setlength{\tabcolsep}{6pt}
%   \begin{tabular}{lcccc}
%     \toprule
%     \textbf{Feature Set} & \textbf{Model} & \textbf{Acc} & \textbf{F1} & \textbf{AUC} \\
%     \midrule
%     SQL                              & RF  & 0.860 & 0.818 & 0.931 \\
%     Static                           & XGB & 0.951 & 0.938 & 0.987 \\
%     Static + Graph                   & XGB & 0.961 & 0.951 & 0.991 \\
%     Static + SQL                     & XGB & 0.965 & 0.956 & 0.989 \\
%     Static + Dynamic                 & XGB & 0.973 & 0.967 & 0.995 \\
%     Static + Graph + SQL             & XGB & 0.969 & 0.961 & 0.992 \\
%     Static + Graph + SQL + Dynamic   & XGB & \textbf{0.980} & \textbf{0.975} & \textbf{0.997} \\
%     \bottomrule
%   \end{tabular}
% \end{table}


% %  Table: Time-optimal method prediction (best_time_method)
% %  Shows best model per feature set. Bold = best in column.

% \begin{table}[t]
%   \centering
%   \caption{%
%     Ablation study for \emph{time-optimal method} prediction
%     (train=6{,}164 / test=1{,}541; class split 85.7\% / 14.3\%).
%     Each row reports the best-performing model for that feature set.
%   }
%   \label{tab:ablation_time}
%   \small
%   \setlength{\tabcolsep}{6pt}
%   \begin{tabular}{lcccc}
%     \toprule
%     \textbf{Feature Set} & \textbf{Model} & \textbf{Acc} & \textbf{F1} & \textbf{AUC} \\
%     \midrule
%     SQL                              & XGB & 0.923 & 0.752 & 0.936 \\
%     Static                           & XGB & 0.920 & 0.738 & 0.955 \\
%     Static + Graph                   & XGB & 0.922 & 0.742 & 0.962 \\
%     Static + SQL                     & XGB & 0.940 & 0.798 & 0.969 \\
%     Static + Dynamic                 & RF  & 0.940 & 0.769 & 0.974 \\
%     Static + Graph + SQL             & XGB & 0.933 & 0.778 & 0.978 \\
%     Static + Graph + SQL + Dynamic   & XGB & \textbf{0.953} & \textbf{0.844} & \textbf{0.986} \\
%     \bottomrule
%   \end{tabular}
% \end{table}


% ----------------------------------------------------------------
%  Requires: booktabs, multicolumn only. No extra packages.
%  table* spans both columns.
%  Acc + F1 at 2dp. Train(s) at 2dp. Infer(ms) at 1dp.
%  Bold = best overall (XGB, full feature set)
%  Underline = second best (RF, full feature set)
% ----------------------------------------------------------------
% --- Table 2: Time-optimal method ---

% \begin{table*}[!ht]
%   \centering
%   \caption{Ablation study for \emph{time-optimal method} prediction
%     (80\% - 20\% train-test split).
%     \textbf{Bold}: best overall result (XGB, full feature set);
%     \underline{underline}: second best (RF, full feature set).
%     Acc: accuracy; F1: macro F1-score; Tr: training time (s); Inf: inference time (ms).}
%   \label{tab:ablation_time}
%   \scriptsize
%   \setlength{\tabcolsep}{3pt}
%   \renewcommand{\arraystretch}{1.05}
%   \begin{tabular}{l cccc cccc cccc cccc cccc}
%     \toprule
%     & \multicolumn{4}{c}{\textbf{LR}}
%     & \multicolumn{4}{c}{\textbf{SVM}}
%     & \multicolumn{4}{c}{\textbf{DT}}
%     & \multicolumn{4}{c}{\textbf{RF}}
%     & \multicolumn{4}{c}{\textbf{XGB}} \\
%     \cmidrule(lr){2-5}\cmidrule(lr){6-9}\cmidrule(lr){10-13}
%     \cmidrule(lr){14-17}\cmidrule(lr){18-21}
%     \textbf{Feature Set}
%       & Acc & F1 & Tr & Inf
%       & Acc & F1 & Tr & Inf
%       & Acc & F1 & Tr & Inf
%       & Acc & F1 & Tr & Inf
%       & Acc & F1 & Tr & Inf \\
%     \midrule
%     SQL
%       & 0.79 & 0.54 & 0.04 & 25.1
%       & 0.80 & 0.57 & 1.70 & 241.9
%       & 0.91 & 0.70 & 0.03 & 10.6
%       & 0.93 & 0.75 & 0.73 & 102.9
%       & 0.92 & 0.75 & 1.41 & 18.6 \\
%     Static
%       & 0.80 & 0.55 & 0.04 & 9.9
%       & 0.80 & 0.55 & 2.10 & 253.6
%       & 0.90 & 0.68 & 0.05 & 10.7
%       & 0.92 & 0.72 & 0.76 & 134.3
%       & 0.92 & 0.74 & 0.40 & 19.0 \\
%     Static + Graph
%       & 0.84 & 0.60 & 0.06 & 10.5
%       & 0.81 & 0.56 & 2.26 & 266.0
%       & 0.90 & 0.66 & 0.08 & 10.6
%       & 0.92 & 0.73 & 0.75 & 134.6
%       & 0.92 & 0.74 & 0.41 & 19.0 \\
%     Static + SQL
%       & 0.86 & 0.64 & 0.07 & 10.3
%       & 0.86 & 0.65 & 2.21 & 229.7
%       & 0.93 & 0.76 & 0.07 & 11.4
%       & 0.94 & 0.78 & 0.76 & 124.6
%       & 0.94 & 0.80 & 0.33 & 12.6 \\
%     Static + Dynamic
%       & 0.86 & 0.66 & 0.04 & 10.0
%       & 0.86 & 0.65 & 1.64 & 209.3
%       & 0.92 & 0.73 & 0.05 & 22.4
%       & 0.94 & 0.77 & 0.72 & 121.7
%       & 0.94 & 0.79 & 0.27 & 13.2 \\
%     Static + Graph + SQL
%       & 0.87 & 0.67 & 0.09 & 12.5
%       & 0.84 & 0.62 & 2.79 & 229.5
%       & 0.92 & 0.74 & 0.11 & 10.2
%       & 0.94 & 0.78 & 0.78 & 137.5
%       & 0.93 & 0.78 & 0.37 & 13.3 \\
%     Static + Graph + SQL + Dyn.
%       & 0.89 & 0.70 & 0.12 & 15.4
%       & 0.88 & 0.69 & 1.80 & 195.1
%       & 0.94 & 0.80 & 0.10 & 10.2
%       & \underline{0.95} & \underline{0.83} & \underline{0.79} & \underline{131.9}
%       & \textbf{0.95} & \textbf{0.84} & \textbf{0.36} & \textbf{13.9} \\
%     \bottomrule
%   \end{tabular}
% \end{table*}
% % --- Table 1: Memory-optimal method ---

% \begin{table*}[!ht]
%   \centering
%   \caption{Ablation study for \emph{memory-optimal method} prediction
%     (80\% - 20\% train-test split).
%     \textbf{Bold}: best overall result (XGB, full feature set);
%     \underline{underline}: second best (RF, full feature set).
%     Acc: accuracy; F1: macro F1-score; Tr: training time (s); Inf: inference time (ms).}
%   \label{tab:ablation_mem}
%   \scriptsize
%   \setlength{\tabcolsep}{3pt}
%   \renewcommand{\arraystretch}{1.05}
%   \begin{tabular}{l cccc cccc cccc cccc cccc}
%     \toprule
%     & \multicolumn{4}{c}{\textbf{LR}}
%     & \multicolumn{4}{c}{\textbf{SVM}}
%     & \multicolumn{4}{c}{\textbf{DT}}
%     & \multicolumn{4}{c}{\textbf{RF}}
%     & \multicolumn{4}{c}{\textbf{XGB}} \\
%     \cmidrule(lr){2-5}\cmidrule(lr){6-9}\cmidrule(lr){10-13}
%     \cmidrule(lr){14-17}\cmidrule(lr){18-21}
%     \textbf{Feature Set}
%       & Acc & F1 & Tr & Inf
%       & Acc & F1 & Tr & Inf
%       & Acc & F1 & Tr & Inf
%       & Acc & F1 & Tr & Inf
%       & Acc & F1 & Tr & Inf \\
%     \midrule
%     SQL
%       & 0.75 & 0.72 & 0.04 & 11.0
%       & 0.75 & 0.72 & 1.53 & 207.0
%       & 0.83 & 0.79 & 0.03 & 7.4
%       & 0.86 & 0.82 & 0.72 & 108.1
%       & 0.85 & 0.81 & 1.36 & 16.2 \\
%     Static
%       & 0.87 & 0.85 & 0.03 & 11.9
%       & 0.87 & 0.84 & 1.06 & 128.3
%       & 0.93 & 0.92 & 0.04 & 7.1
%       & 0.95 & 0.93 & 0.69 & 107.4
%       & 0.95 & 0.94 & 0.27 & 15.8 \\
%     Static + Graph
%       & 0.90 & 0.87 & 0.05 & 11.7
%       & 0.89 & 0.87 & 1.01 & 115.3
%       & 0.93 & 0.91 & 0.08 & 7.7
%       & 0.96 & 0.94 & 0.70 & 93.7
%       & 0.96 & 0.95 & 0.30 & 13.5 \\
%     Static + SQL
%       & 0.89 & 0.87 & 0.06 & 11.3
%       & 0.89 & 0.87 & 0.98 & 111.3
%       & 0.94 & 0.93 & 0.07 & 7.1
%       & 0.96 & 0.95 & 0.70 & 105.9
%       & 0.97 & 0.96 & 0.29 & 12.5 \\
%     Static + Dynamic
%       & 0.94 & 0.92 & 0.03 & 11.8
%       & 0.94 & 0.93 & 0.69 & 94.7
%       & 0.95 & 0.94 & 0.04 & 7.1
%       & 0.97 & 0.96 & 0.68 & 92.6
%       & 0.97 & 0.97 & 0.26 & 12.5 \\
%     Static + Graph + SQL
%       & 0.90 & 0.88 & 0.07 & 10.5
%       & 0.90 & 0.88 & 1.18 & 119.4
%       & 0.94 & 0.92 & 0.10 & 10.1
%       & 0.97 & 0.96 & 0.72 & 118.7
%       & 0.97 & 0.96 & 0.34 & 13.4 \\
%     Static + Graph + SQL + Dyn.
%       & 0.94 & 0.93 & 0.07 & 13.9
%       & 0.94 & 0.93 & 0.90 & 94.3
%       & 0.96 & 0.95 & 0.10 & 9.8
%       & \underline{0.98} & \underline{0.97} & \underline{0.72} & \underline{107.6}
%       & \textbf{0.98} & \textbf{0.98} & \textbf{0.34} & \textbf{13.5} \\
%     \bottomrule
%   \end{tabular}
% \end{table*}
% \begin{table*}[!ht]
%   \centering
%   \caption{Ablation study for \emph{time-optimal method} prediction
%     (80\% - 20\% train-test split).
%     \textbf{Bold}: best overall result (XGB, full feature set);
%     \underline{underline}: second best (RF, full feature set).
%     Acc: accuracy; F1: macro F1-score; Tr: training time (s); Inf: inference time (ms).}
%   \label{tab:ablation_time}
%   \scriptsize
%   \setlength{\tabcolsep}{3pt}
%   \renewcommand{\arraystretch}{1.05}
%   \begin{tabular}{l cccc cccc cccc cccc cccc}
%     \toprule
%     & \multicolumn{4}{c}{\textbf{LR}}
%     & \multicolumn{4}{c}{\textbf{SVM}}
%     & \multicolumn{4}{c}{\textbf{DT}}
%     & \multicolumn{4}{c}{\textbf{RF}}
%     & \multicolumn{4}{c}{\textbf{XGB}} \\
%     \cmidrule(lr){2-5}\cmidrule(lr){6-9}\cmidrule(lr){10-13}
%     \cmidrule(lr){14-17}\cmidrule(lr){18-21}
%     \textbf{Feature Set}
%       & Acc & F1 & Tr & Inf
%       & Acc & F1 & Tr & Inf
%       & Acc & F1 & Tr & Inf
%       & Acc & F1 & Tr & Inf
%       & Acc & F1 & Tr & Inf \\
%     \midrule
%     SQL
%       & 0.79 & 0.54 & \underline{0.04} & 25.1
%       & 0.80 & 0.57 & 1.70 & 241.9
%       & 0.91 & 0.70 & \textbf{0.03} & 10.6
%       & 0.93 & 0.75 & 0.73 & 102.9
%       & 0.92 & 0.75 & 1.41 & 18.6 \\
%     Static
%       & 0.80 & 0.55 & \underline{0.04} & \textbf{9.9}
%       & 0.80 & 0.55 & 2.10 & 253.6
%       & 0.90 & 0.68 & 0.05 & 10.7
%       & 0.92 & 0.72 & 0.76 & 134.3
%       & 0.92 & 0.74 & 0.40 & 19.0 \\
%     Static + Graph
%       & 0.84 & 0.60 & 0.06 & 10.5
%       & 0.81 & 0.56 & 2.26 & 266.0
%       & 0.90 & 0.66 & 0.08 & 10.6
%       & 0.92 & 0.73 & 0.75 & 134.6
%       & 0.92 & 0.74 & 0.41 & 19.0 \\
%     Static + SQL
%       & 0.86 & 0.64 & 0.07 & 10.3
%       & 0.86 & 0.65 & 2.21 & 229.7
%       & 0.93 & 0.76 & 0.07 & 11.4
%       & 0.94 & 0.78 & 0.76 & 124.6
%       & 0.94 & 0.80 & 0.33 & 12.6 \\
%     Static + Dynamic
%       & 0.86 & 0.66 & \underline{0.04} & \underline{10.0}
%       & 0.86 & 0.65 & 1.64 & 209.3
%       & 0.92 & 0.73 & 0.05 & 22.4
%       & 0.94 & 0.77 & 0.72 & 121.7
%       & 0.94 & 0.79 & 0.27 & 13.2 \\
%     Static + Graph + SQL
%       & 0.87 & 0.67 & 0.09 & 12.5
%       & 0.84 & 0.62 & 2.79 & 229.5
%       & 0.92 & 0.74 & 0.11 & 10.2
%       & 0.94 & 0.78 & 0.78 & 137.5
%       & 0.93 & 0.78 & 0.37 & 13.3 \\
%     Static + Graph + SQL + Dyn.
%       & 0.89 & 0.70 & 0.12 & 15.4
%       & 0.88 & 0.69 & 1.80 & 195.1
%       & 0.94 & 0.80 & 0.10 & 10.2
%       & \underline{0.95} & \underline{0.83} & 0.79 & 131.9
%       & \textbf{0.95} & \textbf{0.84} & 0.36 & 13.9 \\
%     \bottomrule
%   \end{tabular}
% \end{table*}
% % --- Table 1: Memory-optimal method ---

% \begin{table*}[!ht]
%   \centering
%   \caption{Ablation study for \emph{memory-optimal method} prediction
%     (80\% - 20\% train-test split).
%     \textbf{Bold}: best overall result (XGB, full feature set);
%     \underline{underline}: second best (RF, full feature set).
%     Acc: accuracy; F1: macro F1-score; Tr: training time (s); Inf: inference time (ms).}
%   \label{tab:ablation_mem}
%   \scriptsize
%   \setlength{\tabcolsep}{3pt}
%   \renewcommand{\arraystretch}{1.05}
%   \begin{tabular}{l cccc cccc cccc cccc cccc}
%     \toprule
%     & \multicolumn{4}{c}{\textbf{LR}}
%     & \multicolumn{4}{c}{\textbf{SVM}}
%     & \multicolumn{4}{c}{\textbf{DT}}
%     & \multicolumn{4}{c}{\textbf{RF}}
%     & \multicolumn{4}{c}{\textbf{XGB}} \\
%     \cmidrule(lr){2-5}\cmidrule(lr){6-9}\cmidrule(lr){10-13}
%     \cmidrule(lr){14-17}\cmidrule(lr){18-21}
%     \textbf{Feature Set}
%       & Acc & F1 & Tr & Inf
%       & Acc & F1 & Tr & Inf
%       & Acc & F1 & Tr & Inf
%       & Acc & F1 & Tr & Inf
%       & Acc & F1 & Tr & Inf \\
%     \midrule
%     SQL
%       & 0.75 & 0.72 & \underline{0.04} & 11.0
%       & 0.75 & 0.72 & 1.53 & 207.0
%       & 0.83 & 0.79 & \textbf{0.03} & \underline{7.4}
%       & 0.86 & 0.82 & 0.72 & 108.1
%       & 0.85 & 0.81 & 1.36 & 16.2 \\
%     Static
%       & 0.87 & 0.85 & \textbf{0.03} & 11.9
%       & 0.87 & 0.84 & 1.06 & 128.3
%       & 0.93 & 0.92 & \underline{0.04} & \textbf{7.1}
%       & 0.95 & 0.93 & 0.69 & 107.4
%       & 0.95 & 0.94 & 0.27 & 15.8 \\
%     Static + Graph
%       & 0.90 & 0.87 & 0.05 & 11.7
%       & 0.89 & 0.87 & 1.01 & 115.3
%       & 0.93 & 0.91 & 0.08 & 7.7
%       & 0.96 & 0.94 & 0.70 & 93.7
%       & 0.96 & 0.95 & 0.30 & 13.5 \\
%     Static + SQL
%       & 0.89 & 0.87 & 0.06 & 11.3
%       & 0.89 & 0.87 & 0.98 & 111.3
%       & 0.94 & 0.93 & 0.07 & \textbf{7.1}
%       & 0.96 & 0.95 & 0.70 & 105.9
%       & 0.97 & 0.96 & 0.29 & 12.5 \\
%     Static + Dynamic
%       & 0.94 & 0.92 & \textbf{0.03} & 11.8
%       & 0.94 & 0.93 & 0.69 & 94.7
%       & 0.95 & 0.94 & \underline{0.04} & \textbf{7.1}
%       & 0.97 & 0.96 & 0.68 & 92.6
%       & 0.97 & 0.97 & 0.26 & 12.5 \\
%     Static + Graph + SQL
%       & 0.90 & 0.88 & 0.07 & 10.5
%       & 0.90 & 0.88 & 1.18 & 119.4
%       & 0.94 & 0.92 & 0.10 & 10.1
%       & 0.97 & 0.96 & 0.72 & 118.7
%       & 0.97 & 0.96 & 0.34 & 13.4 \\
%     Static + Graph + SQL + Dyn.
%       & 0.94 & 0.93 & 0.07 & 13.9
%       & 0.94 & 0.93 & 0.90 & 94.3
%       & 0.96 & 0.95 & 0.10 & 9.8
%       & \underline{0.98} & \underline{0.97} & 0.72 & 107.6
%       & \textbf{0.98} & \textbf{0.98} & 0.34 & 13.5 \\
%     \bottomrule
%   \end{tabular}
% \end{table*}




% --- Memory-optimal method ---

\begin{table*}[!ht]
  \centering
  \caption{Ablation study for \emph{memory-optimal method} prediction
    (80\% - 20\% train-test split).
    \textbf{Bold}: best overall result (XGB, full feature set);
    \underline{underline}: second best (RF, full feature set).
    Acc: accuracy; F1: macro F1-score.}
  \label{tab:ablation_mem}
  \scriptsize
  \setlength{\tabcolsep}{3pt}
  \renewcommand{\arraystretch}{1.05}
  \begin{tabular}{l cc cc cc cc cc}
    \toprule
    & \multicolumn{2}{c}{\textbf{LR}}
    & \multicolumn{2}{c}{\textbf{SVM}}
    & \multicolumn{2}{c}{\textbf{DT}}
    & \multicolumn{2}{c}{\textbf{RF}}
    & \multicolumn{2}{c}{\textbf{XGB}} \\
    \cmidrule(lr){2-3}\cmidrule(lr){4-5}\cmidrule(lr){6-7}
    \cmidrule(lr){8-9}\cmidrule(lr){10-11}
    \textbf{Feature Set}
      & Acc & F1
      & Acc & F1
      & Acc & F1
      & Acc & F1
      & Acc & F1 \\
    \midrule
    SQL
      & 0.75 & 0.72
      & 0.75 & 0.72
      & 0.83 & 0.79
      & 0.86 & 0.82
      & 0.85 & 0.81 \\
    Static
      & 0.87 & 0.85
      & 0.87 & 0.84
      & 0.93 & 0.92
      & 0.95 & 0.93
      & 0.95 & 0.94 \\
    Static + Graph
      & 0.90 & 0.87
      & 0.89 & 0.87
      & 0.93 & 0.91
      & 0.96 & 0.94
      & 0.96 & 0.95 \\
    Static + SQL
      & 0.89 & 0.87
      & 0.89 & 0.87
      & 0.94 & 0.93
      & 0.96 & 0.95
      & 0.97 & 0.96 \\
    Static + Dynamic
      & 0.94 & 0.92
      & 0.94 & 0.93
      & 0.95 & 0.94
      & 0.97 & 0.96
      & 0.97 & 0.97 \\
    Static + Graph + SQL
      & 0.90 & 0.88
      & 0.90 & 0.88
      & 0.94 & 0.92
      & 0.97 & 0.96
      & 0.97 & 0.96 \\
    Static + Graph + SQL + Dyn.
      & 0.94 & 0.93
      & 0.94 & 0.93
      & 0.96 & 0.95
      & \underline{0.98} & \underline{0.97}
      & \textbf{0.98} & \textbf{0.98} \\
    \bottomrule
  \end{tabular}
\end{table*}


We conduct a feature-group ablation study on the 7705 dataset. 
Tables~\ref{tab:ablation_time} and~\ref{tab:ablation_mem} report the performance of five models (Logistic Regression, Linear SVM, Decision Tree, Random Forest, and XGBoost) for runtime- and memory-optimal routing, respectively. 
% XGBoost consistently achieves the best performance in terms of both accuracy and F1 score, motivating its use in our feature-group study in Section~\ref{sec:eva}. 
We observe that combining all feature groups yields the best overall performance. 
In some cases, such as Logistic Regression in Table~\ref{tab:ablation_mem}, adding dynamic features also leads to a noticeable improvement.



 

% \section{Training and Inference Time}

% This experiment further confirms that XGBoost is well-suited to our routing task. 
% In both tables, XGBoost achieves substantially lower inference time than Random Forest, the second strongest baseline in terms of predictive performance. 
% Although Decision Trees have slightly lower inference time than XGBoost, XGBoost achieves higher accuracy and F1 score, providing a better trade-off between prediction quality and inference efficiency.

% \begin{table*}[!ht]
%   \centering
%   \caption{Training and inference time for \emph{time-optimal method} prediction
%     (80\% - 20\% train-test split).
%     \textbf{Bold}: lowest value;
%     \underline{underline}: second-lowest value.
%     Tr: training time (s); Inf: inference time (ms).}
%   \label{tab:ablation_time_cost}
%   \scriptsize
%   \setlength{\tabcolsep}{3pt}
%   \renewcommand{\arraystretch}{1.05}
%   \begin{tabular}{l cc cc cc cc cc}
%     \toprule
%     & \multicolumn{2}{c}{\textbf{LR}}
%     & \multicolumn{2}{c}{\textbf{SVM}}
%     & \multicolumn{2}{c}{\textbf{DT}}
%     & \multicolumn{2}{c}{\textbf{RF}}
%     & \multicolumn{2}{c}{\textbf{XGB}} \\
%     \cmidrule(lr){2-3}\cmidrule(lr){4-5}\cmidrule(lr){6-7}
%     \cmidrule(lr){8-9}\cmidrule(lr){10-11}
%     \textbf{Feature Set}
%       & Tr & Inf
%       & Tr & Inf
%       & Tr & Inf
%       & Tr & Inf
%       & Tr & Inf \\
%     \midrule
%     SQL
%       & \underline{0.04} & 25.1
%       & 1.70 & 241.9
%       & \textbf{0.03} & 10.6
%       & 0.73 & 102.9
%       & 1.41 & 18.6 \\
%     Static
%       & \underline{0.04} & \textbf{9.9}
%       & 2.10 & 253.6
%       & 0.05 & 10.7
%       & 0.76 & 134.3
%       & 0.40 & 19.0 \\
%     Static + Graph
%       & 0.06 & 10.5
%       & 2.26 & 266.0
%       & 0.08 & 10.6
%       & 0.75 & 134.6
%       & 0.41 & 19.0 \\
%     Static + SQL
%       & 0.07 & 10.3
%       & 2.21 & 229.7
%       & 0.07 & 11.4
%       & 0.76 & 124.6
%       & 0.33 & 12.6 \\
%     Static + Dynamic
%       & \underline{0.04} & \underline{10.0}
%       & 1.64 & 209.3
%       & 0.05 & 22.4
%       & 0.72 & 121.7
%       & 0.27 & 13.2 \\
%     Static + Graph + SQL
%       & 0.09 & 12.5
%       & 2.79 & 229.5
%       & 0.11 & 10.2
%       & 0.78 & 137.5
%       & 0.37 & 13.3 \\
%     Static + Graph + SQL + Dyn.
%       & 0.12 & 15.4
%       & 1.80 & 195.1
%       & 0.10 & 10.2
%       & 0.79 & 131.9
%       & 0.36 & 13.9 \\
%     \bottomrule
%   \end{tabular}
% \end{table*}

% \begin{table*}[!ht]
%   \centering
%   \caption{Training and inference time for \emph{memory-optimal method} prediction
%     (80\% - 20\% train-test split).
%     \textbf{Bold}: lowest value;
%     \underline{underline}: second-lowest value.
%     Tr: training time (s); Inf: inference time (ms).}
%   \label{tab:ablation_mem_cost}
%   \scriptsize
%   \setlength{\tabcolsep}{3pt}
%   \renewcommand{\arraystretch}{1.05}
%   \begin{tabular}{l cc cc cc cc cc}
%     \toprule
%     & \multicolumn{2}{c}{\textbf{LR}}
%     & \multicolumn{2}{c}{\textbf{SVM}}
%     & \multicolumn{2}{c}{\textbf{DT}}
%     & \multicolumn{2}{c}{\textbf{RF}}
%     & \multicolumn{2}{c}{\textbf{XGB}} \\
%     \cmidrule(lr){2-3}\cmidrule(lr){4-5}\cmidrule(lr){6-7}
%     \cmidrule(lr){8-9}\cmidrule(lr){10-11}
%     \textbf{Feature Set}
%       & Tr & Inf
%       & Tr & Inf
%       & Tr & Inf
%       & Tr & Inf
%       & Tr & Inf \\
%     \midrule
%     SQL
%       & \underline{0.04} & 11.0
%       & 1.53 & 207.0
%       & \textbf{0.03} & \underline{7.4}
%       & 0.72 & 108.1
%       & 1.36 & 16.2 \\
%     Static
%       & \textbf{0.03} & 11.9
%       & 1.06 & 128.3
%       & \underline{0.04} & \textbf{7.1}
%       & 0.69 & 107.4
%       & 0.27 & 15.8 \\
%     Static + Graph
%       & 0.05 & 11.7
%       & 1.01 & 115.3
%       & 0.08 & 7.7
%       & 0.70 & 93.7
%       & 0.30 & 13.5 \\
%     Static + SQL
%       & 0.06 & 11.3
%       & 0.98 & 111.3
%       & 0.07 & \textbf{7.1}
%       & 0.70 & 105.9
%       & 0.29 & 12.5 \\
%     Static + Dynamic
%       & \textbf{0.03} & 11.8
%       & 0.69 & 94.7
%       & \underline{0.04} & \textbf{7.1}
%       & 0.68 & 92.6
%       & 0.26 & 12.5 \\
%     Static + Graph + SQL
%       & 0.07 & 10.5
%       & 1.18 & 119.4
%       & 0.10 & 10.1
%       & 0.72 & 118.7
%       & 0.34 & 13.4 \\
%     Static + Graph + SQL + Dyn.
%       & 0.07 & 13.9
%       & 0.90 & 94.3
%       & 0.10 & 9.8
%       & 0.72 & 107.6
%       & 0.34 & 13.5 \\
%     \bottomrule
%   \end{tabular}
% \end{table*}
 

\section{Training and Inference Time}
\label{sec:traintime}
We report the computational cost of each classifier under the all-features setting.
Table~\ref{tab:ablation_cost_all_features} reports training and inference time separately for runtime- and memory-optimal routing, with all values in seconds.
Training time is measured over 6,164 training samples, whereas inference time is measured over 1,541 test samples; therefore, the two measurements should be interpreted independently rather than compared directly.

Combining these results with the accuracy and F1 results in Tables~\ref{tab:ablation_time} and~\ref{tab:ablation_mem}, we observe that XGBoost provides the best overall trade-off between predictive performance and inference efficiency.
Although Decision Trees have the lowest inference time for both routing tasks, requiring 0.0102 s for runtime-optimal routing and 0.0098 s for memory-optimal routing, XGBoost achieves substantially higher accuracy and F1 score while remaining highly efficient at inference time.
Compared with Random Forest, the second strongest classifier in predictive performance, XGBoost is considerably faster during inference: 0.0139 s versus 0.1319 s for runtime-optimal routing and 0.0135 s versus 0.1076 s for memory-optimal routing, measured over the same 1,541 inference samples.
XGBoost also maintains moderate training cost, requiring 0.36 s and 0.34 s for runtime- and memory-optimal routing, respectively.
These results further support our use of XGBoost as the routing model.

% \begin{table}[h]
%   \centering
%   \caption{Training and inference time under the all-features setting
%     (80\%--20\% train-test split).
%     \textbf{Bold}: lowest value;
%     \underline{underline}: second-lowest value.
%     Tr: training time (s); Inf: inference time (ms).}
%   \label{tab:ablation_cost_all_features}
%   \scriptsize
%   \setlength{\tabcolsep}{4pt}
%   \renewcommand{\arraystretch}{1.05}
%   \begin{tabular}{l cc cc}
%     \toprule
%     & \multicolumn{2}{c}{\textbf{Runtime-Optimal Routing}}
%     & \multicolumn{2}{c}{\textbf{Memory-Optimal Routing}} \\
%     \cmidrule(lr){2-3}\cmidrule(lr){4-5}
%     \textbf{Model}
%       & Tr & Inf
%       & Tr & Inf \\
%     \midrule
%     LR
%       & \underline{0.12} & 15.4
%       & \textbf{0.07} & 13.9 \\
%     SVM
%       & 1.80 & 195.1
%       & 0.90 & 94.3 \\
%     DT
%       & \textbf{0.10} & \textbf{10.2}
%       & \underline{0.10} & \textbf{9.8} \\
%     RF
%       & 0.79 & 131.9
%       & 0.72 & 107.6 \\
%     XGB
%       & 0.36 & \underline{13.9}
%       & 0.34 & \underline{13.5} \\
%     \bottomrule
%   \end{tabular}
% \end{table}
\begin{table}[thb]
  \centering
  \caption{Training and inference time under the all-features setting
    using an 80\%--20\% train-test split.
    All values are total elapsed time in seconds.
    Training was measured over 6,164 samples, while inference was measured
    over 1,541 samples.
    \textbf{Bold}: lowest value;
    \underline{underline}: second-lowest value within each section and routing setting.}
  \label{tab:ablation_cost_all_features}
  \scriptsize
  \setlength{\tabcolsep}{5pt}
  \renewcommand{\arraystretch}{1.08}
  \begin{tabular}{lcc}
    \toprule
    \multicolumn{3}{c}{\textbf{Training time over 6,164 samples (s)}} \\
    \midrule
    \textbf{Model}
      & \textbf{Runtime-Optimal Routing}
      & \textbf{Memory-Optimal Routing} \\
    \midrule
    LR
      & \underline{0.12}
      & \textbf{0.07} \\
    SVM
      & 1.80
      & 0.90 \\
    DT
      & \textbf{0.10}
      & \underline{0.10} \\
    RF
      & 0.79
      & 0.72 \\
    XGB
      & 0.36
      & 0.34 \\
    \midrule
    \multicolumn{3}{c}{\textbf{Inference time over 1,541 samples (s)}} \\
    \midrule
    \textbf{Model}
      & \textbf{Runtime-Optimal Routing}
      & \textbf{Memory-Optimal Routing} \\
    \midrule
    LR
      & 0.0154
      & 0.0139 \\
    SVM
      & 0.1951
      & 0.0943 \\
    DT
      & \textbf{0.0102}
      & \textbf{0.0098} \\
    RF
      & 0.1319
      & 0.1076 \\
    XGB
      & \underline{0.0139}
      & \underline{0.0135} \\
    \bottomrule
  \end{tabular}
\end{table}
\section{Tensor-Aware DBMSs}
\label{app:tensor-aware-dbms}

\begin{figure*}[t]
    \centering
    \includegraphics[width=\linewidth]{images/scidb_on_box_pc.pdf}
    \caption{Additional DB comparison on 11 small \sys circuits with varying output density. 
    Top row: runtime; bottom row: peak memory. The plots group circuits by qubit count,
    output-density bin, and gate count.}
    \label{fig:addscidbcomp}
\end{figure*}

Beyond the RDBMS engines used in the main evaluation, tensor-aware DBMSs
such as TileDB and SciDB are natural candidates for quantum circuit simulation.
They expose array-oriented storage and execution abstractions that appear well matched
to tensor contractions. However, using them fairly requires significant effort because \sys currently
emits SQL workloads whose core operations are joins and aggregations.

% \para{TileDB.}
TileDB is an array storage engine for dense and sparse multidimensional arrays.\footnote{
TileDB documentation: \url{https://docs.tiledb.com/main/}. TileDB-Py API:
\url{https://tiledb-inc-tiledb.readthedocs-hosted.com/projects/tiledb-py/en/stable/python-api.html}.}
A direct TileDB evaluation would require an array-native implementation of the simulator:
one would need to choose array schemas, chunk sizes, qubit-index layouts, and contraction
kernels, rather than executing the same SQL workload used by PostgreSQL, SQLite, and
DuckDB. The closest SQL-oriented possibility, \texttt{TileDB-MariaDB/MyTile}, is deprecated and supports
only limited pushdown, which makes it unsuitable as a fair drop-in replacement for our
join-and-aggregate SQL workload.\footnote{
TileDB-MariaDB/MyTile documentation:
\url{https://github.com/TileDB-Inc/TileDB-MariaDB}.}
We therefore do not include TileDB in the timing comparison. We view a TileDB-native
implementation as important future work rather than a direct backend substitution.

% \para{SciDB.}
SciDB is closer to our setting because it provides array operators and a query interface
for composing them.\footnote{
SciDB-Py documentation: \url{https://paradigm4.github.io/SciDB-Py/guide.html}.}
Since \sys emits SQL, we made a preliminary SciDB implementation by translating each tensor
contraction into SciDB's Array Functional Language (AFL)-style array operations. In principle, an entire circuit could be
encoded as one nested array expression. In practice, the generated expressions become
large and fragile for full circuits, so our prototype splits the simulation at tensor
contraction boundaries and materializes intermediate arrays. This makes the prototype
robust, but it removes some global optimization opportunities available to RDBMSs when
they optimize a CTE chain. Thus, please consider the results below as preliminary
results rather than as results from a fully optimized SciDB implementation.

\medskip
\para{Experimental setting}
Figure~\ref{fig:addscidbcomp} reports a comparison on 11 \sys circuits \footnote{\url{https://github.com/InfiniData-Lab/InferQ/blob/main/analysis/sample_csvs/small_sample_scidb.csv}} with
3--7 qubits and 4--17 gates, spanning sparse, mixed, and dense output-density bins.
We compare PostgreSQL, DuckDB, SQLite, Umbra, SciDB, and two NumPy baselines:
\texttt{np-one-shot}, an exact tensor-contraction baseline, and \texttt{np-mps}, an
MPS-based approximate baseline. 
% The experiment is CPU-only and was run on a local
% 32GB-RAM machine; PostgreSQL, SciDB, and Umbra run in Docker containers, while SQLite
% and DuckDB use Python packages. Since this machine differs from the main experimental
% platform, this experiment is intended as a qualitative engine discussion rather than a
% replacement for the main evaluation.

\medskip
\para{Results}
We have two main observations from Figure~\ref{fig:addscidbcomp}. First, among DB-backed
methods, SQLite remains the strongest backend. It beats
\texttt{np-one-shot} in peak memory for all 11 circuits and in runtime for 4 of the
11 circuits. PostgreSQL and Umbra remain competitive in memory but are slower than
SQLite. Second, the direct SciDB prototype has substantial overhead: it is consistently
slower and uses more memory than the RDBMS backends on these small circuits. The
\texttt{np-mps} baseline is often fast, but it is approximate, whereas the DB-backed
methods and \texttt{np-one-shot} are exact strong simulators.

\medskip
  \vspace{0.2cm}
\noindent
\setlength{\fboxsep}{6pt}
\fcolorbox{black!15}{black!4}{%
  \parbox{\dimexpr\linewidth-2\fboxsep-2\fboxrule\relax}{
\para{Take-away}
Tensor-aware DBMSs are promising, but they are not plug-in replacements for the SQL
workload studied in this paper. To make TileDB or SciDB competitive for quantum circuit
simulation, future work should co-design the simulator with the array engine, including:
(i) qubit-index layout and chunking, (ii) sparse-versus-dense array selection,
(iii) contraction fusion to avoid per-step materialization, and (iv) cost models for
array-native execution. It is an interesting future direction to explore which types of simulation workloads 
array DBMSs would be most promising for.
  }}
%   for larger workloads
% where the overhead can be amortized, especially circuits whose states remain sparse or
% whose tensor contractions have regular chunkable structure. Small circuits and highly
% irregular contraction sequences are less favorable without such array-native optimization.










% \section{Discussion of additional vector databases (SciDB)}
% \begin{figure*}
%     \centering
%     \includegraphics[width=\linewidth]{images/scdb_trends.pdf}
%     \caption{Additional DB comparisons for 11 small InferQ circuits with varying sparsity. Plots on the top specify memory and the bottom time usage}
%     \label{fig:addscidbcomp}
% \end{figure*}


% Apart from SQLite, PostgreSQL and DuckDB, we also discuss the use of vector databases like SciDB,  TileDB and others like UmbraDB.

% InferQ circuits can be converted to SQL for PostgreSQL, DuckDB and SQLite, we add support for UmbraDB and SciDB too. For TileDB, this is is much not convoluted, for running through MariaDB and cloud, hence we omit it from the comparison. For SciDB, we need to adapt our SQL queries to AFL queries. These quantum circuit queries unfortunately are too large for the SciDB single query AFL. Consequently, we need to resort to splitting the query into per tensor contraction AFLs, reducing any join order optimisation benefit that the other DBs offer. Another limitation is that SciDb is also maintained less to the more common PostgreSQL and SQLite engines. 

% We ran a small experiment on 11 small quantum circuits for InferQ on our local laptop machine due to occupancy of the main machine used in other experiments. The specifications are Asus ROG Zephyrus G14 32Gb RAM, AMD Ryzen 9 CPU and Nvidia GeForce RTZ 5070 GPU, which is not used. \rh{Do I need this specification here?}. UmbraDB, SciDB and PostgreSQL were run with docker containers while SQLite and DuckDB have python packages. The experiment serves the purpose to learn more about SciDB for quantum circuits and expand the discussion about database engines. 


% Figure \ref{fig:addscidbcomp} displays the performance for time and memory usage to simulate a sample of 11 InferQ circuits for PostgreSQL, DuckDB, SQLite, SciDB and UmbraDB. Out of the 11 circuits, only SQLite is able to beat its numpy tensor contraction counterparts in 4 / 11 cases for runtime and in all 11 /11 cases for memory optimality. Between the RDBMS systems, we see that for these small InferQ quantum circuits SQLite is consistently the best. SciDB is a couple of order fo magnitude worse to other engines, apart from DuckDB which has an even more expensive run for a circuit with 7 qubits. The numpy MPS tensor network (np-mps) performs well, but provides an approximation of the final output unlike the other systems which are exact strong simulators. 

% Despite UmbraDB falling short to PostgreSQL in time, it is slightly better for memory with both these engines remaining competitive against their exact numpy tensor (np-one-shot) simulator. The ranking amongst these systems remain similar across number of qubits, number of gates and the output sparsity of the quantum circuit. 

% This experiment profiles various RDBMS engines, including SciDB. It is outperformed consistently by the other 4 engines on smaller quantum circuits in InferQ. yet, the experiment highlights that RDBMS simulators can be competitive and suggests that SQLite might be one of the best. 
\section{Support for Existing Benchmarks}
\label{app:benchmark-suites}

In addition to circuits generated by \sys, the evaluation pipeline can ingest circuits
from established benchmark suites: SupermarQ~\cite{tomesh2022supermarqscalablequantumbenchmark},
MQT Bench~\cite{Quetschlich_2023}, and QASM
Bench~\cite{li2022qasmbenchlowlevelqasmbenchmark}. These suites provide workloads that complement our generated corpus. The compositional circuit simulation workload introduced in Sec.~\ref{sec:benchmark}
remains necessary for controlled simulation scale, feature-space coverage, and reproducible
randomized workload generation, while external suites provide an additional validation
source based on named benchmark circuits.

\begin{table}[t]
  \centering
  \caption{External benchmark-suite circuit entries supported by the \sys ingestion pipeline. Per-suite counts are before cross-suite deduplication.}
  \label{tab:benchmark-suite-ingestion}
  \footnotesize
  \begin{tabular}{@{}lr@{}}
    \toprule
    \textbf{Benchmark suite} & \textbf{\# circuit entries} \\
    \midrule
    SupermarQ & 8 \\
    MQT Bench & 34 \\
    QASM Bench & 67 \\
    \midrule
    Total (after deduplication) & 85 \\
    \bottomrule
  \end{tabular}
\end{table}

All imported circuits are normalized to Qiskit \texttt{QuantumCircuit} objects and
tagged with their source suite and benchmark name. The ingestion pipeline then applies
the same SQL generation, feature extraction, and the rest of the workflow used for \sys-generated
circuits. As a result, external-suite circuits can be queried, filtered, and evaluated
together with generated circuits under the same RDBMS, simulator-selection, and
out-of-core workflows.

The per-benchmark counts in Table~\ref{tab:benchmark-suite-ingestion} are before deduplication. Since
the three benchmark suites contain overlapping algorithm families and problem-size variants, \sys
also records a deduplicated inventory; after deduplication, the external benchmark
inventory covers 85 unique algorithm families. The complete per-suite circuit-family
list and ingestion commands are provided in the \emph{Benchmark Algorithms} table in the online documentation.\footnote{\url{https://github.com/InfiniData-Lab/InferQ/blob/main/scripts/benchmark_suites/README.md}}

 
% \section{Deep Profiling of SQLite and Qiskit Aer}
% \label{sec:deep-profile-sqlite-aer}

% % The aggregate evaluation shows that SQLite can be competitive with Qiskit Aer, especially for peak memory, but the aggregate win rates do not explain where this advantage comes from. We have generated and simulated 106 circuits, and profiled the SQLite and Qiskit Aer\footnote{These 106 circuits are sampled from the \num{7705} circuits used in the Sec. 6.}. The goal is to identify the mechanism behind the crossover: when SQLite is faster, whether this is due to sparse relational representation and low overhead; and when Aer is faster, whether SQLite is dominated by SQL execution-plan cost.

% To understand why SQLite uses less peak memory than
% Qiskit Aer on many circuits, we have generated and simulated 106 circuits, and profiled the SQLite and Qiskit Aer\footnote{These 106 circuits are sampled from the \num{7705} circuits used in the Sec. 6.}.
% % This analysis separates SQLite's memory advantage from possible simulator memory
% % overheads and identifies when SQLite's cost is dominated by SQL execution.

% For SQLite, each circuit is compiled into a tensor-contraction program represented as a chain of Common Table Expressions (CTEs). Each contraction step is implemented as a join followed by an aggregation. We record the number of total and contraction CTEs, estimated intermediate-relation sizes, result cardinalities, SQLite execution-plan operators, temporary B-tree usage, automatic-index usage, page-cache overflow counters, peak resident set size (RSS), and wall-clock runtime. For Qiskit Aer, we compare against the fastest successful Aer method for the same circuit, rather than fixing a single simulator representation across all circuits.

% The profiling results in Table~\ref{tab:sqlite-aer-profile} help us understand \emph{when} each engine wins. Across 106 circuits, SQLite is faster in 19 cases, while Qiskit Aer is faster in 87 cases. SQLite wins when the relational workload remains small: the median SQLite-winning circuit has 37 total CTEs, 32 contraction CTEs, a largest intermediate relation of 256\,B, and only 16 final result rows. In these cases, SQLite runs in a median of 0.85\,ms, compared with 2.01\,ms for the fastest successful Aer method. By contrast, when Aer wins, the median workload is larger: 55 total CTEs, 49 contraction CTEs, a largest intermediate relation of 1.0\,KB, and 128 final result rows. In these cases, SQLite takes a median of 1.67\,ms, while Aer takes 0.92\,ms.

% \begin{table}[t]
% \centering
% \small
% \caption{Median properties of completed SQLite--Aer profiling runs, grouped by the faster runtime backend. CTE sizes are estimated from intermediate-relation cardinalities and row widths.}
% \label{tab:sqlite-aer-profile}
% \begin{tabular}{lrr}
% \toprule
%  & SQLite faster & Aer faster \\
% \midrule
% Number of cases & 19 & 87 \\
% Total CTEs & 37 & 55 \\
% Contraction CTEs & 32 & 49 \\
% Total estimated CTE size & 3.0\,KB & 8.6\,KB \\
% Largest intermediate CTE & 256\,B & 1.0\,KB \\
% Final result rows & 16 & 128 \\
% SQLite runtime & 0.85\,ms & 1.67\,ms \\
% Fastest Aer runtime & 2.01\,ms & 0.92\,ms \\
% \bottomrule
% \end{tabular}
% \end{table}

% This result shows that SQLite's advantage is not determined by qubit count alone. Some SQLite-winning circuits have more qubits than Aer-winning circuits, but their relational intermediates remain small. The decisive factor is the size and structure of the induced SQL workload. When the tensor-contraction sequence preserves sparsity, SQLite only materializes a small number of non-zero amplitudes, and the join-and-aggregate pipeline can finish before Aer method setup and simulation overhead dominate. For such circuits, the sparse SQL representation is a good representation of the quantum state in the circuits.

% The same representation becomes unfavorable once the contraction pipeline grows. Longer CTE chains introduce repeated joins, aggregations, and materialization points. SQLite query plans repeatedly use temporary B-trees for grouping and automatic indexes for joins. These are standard relational execution mechanisms, but in a long contraction sequence their cumulative cost dominates the arithmetic. Thus, when Aer wins, it is not because SQLite fails to represent the computation; rather, the relational execution machinery becomes the bottleneck.

% % The timeout cases follow the same pattern. They occur in the high-cost SQL regime, with larger join-and-aggregation structure: the median timeout case has 15 qubits, 99 \texttt{GROUP BY} clauses, and 128 join/equality predicates. These cases indicate that the hard boundary for SQLite is not simply state-vector dimension, but the interaction between output density, intermediate cardinality, CTE-chain length, and the physical operators selected by SQLite. Dense or highly connected circuits can rapidly turn the SQL representation into a long sequence of expensive relational operators.

% The memory measurement results are different from the runtime results. Among 106 circuits, SQLite uses less peak RSS than the best Aer memory configuration in 105 out of 106 cases. Median peak RSS is 167\,MB for SQLite and 171\,MB for Aer. This difference is small in absolute terms for these circuits, but it is consistent: SQLite benefits from storing sparse relational intermediates rather than allocating the simulator representation chosen by Aer. At the same time, this memory advantage does not automatically translate into runtime advantage, because SQLite still pays the cost of repeated SQL planning and execution operators.
% %
% The SQLite profiling results are similar. The page-cache overflow counter is zero in all completed runs, and we do not observe an external-sort. Therefore, the slowdown is not explained by SQLite spilling large page-cache contents to disk. Instead, it might be because the repeated use of temporary B-trees and automatic indexes inside the query plan. 
% % In this setting, temporary B-trees should be interpreted primarily as relational execution structures used for grouping and ordering, not as evidence of out-of-core execution. 
% The dominant cost is the repeated construction and consumption of these structures across the CTE chain.

% Different from RDBMSs, Aer's performance depends strongly on which simulator method is valid and efficient for the circuit. In the profiling runs, the fastest successful Aer method varies across circuits: \texttt{unitary} is fastest in 47 cases, \texttt{automatic} in 40 cases, \texttt{extended\_stabilizer} in 12 cases, and \texttt{statevector} in 7 cases. This confirms that a single fixed Aer method is not a stable baseline for heterogeneous circuits. Some circuits admit efficient specialized simulation, while others require a more general representation. For peak memory, the best Aer configuration is almost always \texttt{statevector}, which is the lowest-RSS Aer method in 105 completed cases.

% Overall, the profiling study explains the SQLite--Aer crossover. SQLite wins when three conditions hold simultaneously: the circuit state remains sparse, the intermediate relations stay small, and the CTE chain is short enough that SQL execution overhead is negligible. Aer wins once the contraction sequence becomes long enough for repeated joins, aggregations, temporary B-trees, and automatic indexes to dominate. The comparison is therefore governed by both data representation and execution strategy: SQLite provides a compact sparse representation for some circuits, while Aer benefits from specialized simulator methods that avoid the overhead of executing a long relational contraction pipeline.

\section{Profiling of SQLite and Qiskit Aer}
\label{sec:deep-profile-sqlite-aer}

To understand why SQLite uses less peak memory than Qiskit Aer on many circuits,
we profile both systems on 162 generated circuits.\footnote{We have sampled these 162 circuits
  from the large \num{7705} circuits used in Section~6. This profiling comparison is separate from the all-backend tuned
winner count reported in Figure~7 and Appendix~G.}

For SQLite, each circuit is compiled into a tensor-contraction program represented
as a chain of Common Table Expressions (CTEs). Each contraction step is implemented
as a join followed by an aggregation. We record the number of total and contraction
CTEs, estimated intermediate-relation sizes, result cardinalities, SQLite
execution-plan operators, temporary B-tree usage, automatic-index usage,
page-cache overflow counters, peak resident set size (RSS), and wall-clock
runtime. For Qiskit Aer, we compare against the fastest   Aer method for
runtime, rather than fixing a single simulator representation across all circuits.
For memory, we compare against the lowest-RSS   Aer configuration for
the same circuit.

The profiling results\footnote{\url{https://github.com/InfiniData-Lab/InferQ/blob/main/analysis/memory_profiling/sqlite_vs_qiskit_summary.csv}} in Table~\ref{tab:sqlite-aer-profile} characterize when
each system is faster. Across 162 circuits, SQLite is faster in 46 cases, while
Qiskit Aer is faster in 116 cases. SQLite is faster when the relational workload
remains small: the median SQLite-winning circuit has 26.5 total CTEs, 22
contraction CTEs, the largest intermediate relation of 256\,B, and 16 final result
rows. In these cases, SQLite runs in a median of 2.03\,ms, compared with
2.82\,ms for the fastest   Aer method. In contrast, when Aer is faster,
the median workload is larger: 55 total CTEs, 49 contraction CTEs, a largest
intermediate relation of 1.0\,KB, and 128 final result rows. In these cases,
SQLite takes a median of 6.63\,ms, while Aer takes 3.09\,ms.

\begin{table}[t]
\centering
\small
\caption{Median properties of completed SQLite--Aer profiling runs, grouped by
the backend with lower runtime. CTE sizes are estimated from intermediate-relation
cardinalities and row widths.}
\label{tab:sqlite-aer-profile}
\begin{tabular}{lrr}
\toprule
 & SQLite faster & Aer faster \\
\midrule
Number of cases & 46 & 116 \\
Total CTEs & 26.5 & 55 \\
Contraction CTEs & 22 & 49 \\
Total estimated CTE size & 2.9\,KB & 10.5\,KB \\
Largest intermediate CTE & 256\,B & 1.0\,KB \\
Final result rows & 16 & 128 \\
SQLite runtime & 2.03\,ms & 6.63\,ms \\
Fastest Aer runtime & 2.82\,ms & 3.09\,ms \\
\bottomrule
\end{tabular}
\end{table}

These results show that SQLite's advantage is not determined by qubit count
alone. Some SQLite-winning circuits have more qubits than Aer-winning circuits,
but their relational intermediates remain small. The decisive factor is the size
and structure of the induced SQL workload. When the tensor-contraction sequence
preserves sparsity, SQLite materializes only a small number of nonzero amplitudes,
and the join-and-aggregate pipeline can complete before Aer's method setup and
simulation overheads dominate. For such circuits, the sparse SQL representation
provides a compact representation of the circuit state.

The same representation becomes less favorable once the contraction pipeline
grows. Longer CTE chains introduce repeated joins, aggregations, and intermediate
execution structures. SQLite query plans use temporary B-trees for grouping in
all completed runs and automatic indexes for joins in 154 out of 162 runs. These
are standard relational execution mechanisms, but in a long contraction sequence
their cumulative cost can dominate the arithmetic. Thus, when Aer is faster, it
is not because SQLite fails to represent the computation; rather, the relational
execution machinery becomes the bottleneck.

The memory measurements follow a different pattern from the runtime measurements.
Among the 162 circuits, SQLite uses less peak RSS than the lowest-RSS  
Aer method in 139 cases. Median peak RSS is 203\,MB for SQLite and 206\,MB for
Aer. This difference is small in absolute terms for these circuits, but it is
consistent with SQLite benefiting from storing sparse relational intermediates
rather than allocating the simulator representation used by Aer. At the same
time, this memory advantage does not automatically translate into a runtime
advantage, because SQLite still pays the cost of repeated SQL execution.

The SQLite-specific profiling counters support this interpretation. The
page-cache overflow counter is zero in all completed runs, and we do not observe
external-sort behavior. Thus, the slowdown is not explained by observed
out-of-core execution in these runs. Instead, the dominant cost is the repeated
construction and consumption of temporary B-trees and automatic indexes across
the CTE chain.

Unlike the RDBMS backends, Aer's performance depends strongly on which simulator
method is applicable and efficient for each circuit. In the profiling runs, the
fastest  Aer method varies across circuits: \texttt{unitary} is fastest
in 46 cases, \texttt{automatic} in 42 cases, \feat{statevector} in 37 cases,
\feat{extended\_stabilizer} in 26 cases, \feat{matrix\_product\_state} in
5 cases, \feat{density\_matrix} in 4 cases, and \feat{stabilizer} in 2
cases. This indicates that a single fixed Aer method is not a stable baseline for
heterogeneous circuits. Some circuits admit efficient specialized simulation,
while others require a more general representation. For peak memory, the
lowest-RSS Aer configuration is usually \feat{statevector} or
\feat{density\_matrix}:  \feat{statevector} has the lowest RSS in 88 cases,
\feat{density\_matrix} in 49 cases, \feat{automatic} in 14 cases,
\feat{matrix\_product\_state} in 9 cases, and \feat{unitary} in 2 cases.

We do not enable Aer cache blocking in this comparison. Our Aer baseline uses the
standard CPU simulator methods, whereas Aer's cache-blocking options (e.g.,
\texttt{blocking\_enable} and \texttt{blocking\_qubits}) are documented for
distributed GPU and/or multi-node simulation~\cite{QiskitAerMultiGPU}. They
partition the state into chunks to reduce data exchange across distributed memory
spaces, rather than providing a disk-backed CPU cache comparable to an RDBMS
buffer manager. Enabling these options would therefore change the simulator
configuration instead of providing a fair comparison to SQLite's
relational execution.  

\medskip
  \vspace{0.2cm}
\noindent
\setlength{\fboxsep}{6pt}
\fcolorbox{black!15}{black!4}{%
  \parbox{\dimexpr\linewidth-2\fboxsep-2\fboxrule\relax}{
\noindent\textbf{Takeaway.}
The relative performance of SQLite and Qiskit Aer is governed by both
representation and execution strategy. SQLite is memory-efficient when sparsity
keeps intermediate relations small, but Aer is faster once repeated joins,
aggregations, temporary B-trees, and automatic indexes make SQL execution
overhead dominate the contraction pipeline.
}} 

% \section{Additional benchmark-suite circuits}
% \label{app:benchmark-suites}

% In addition to the randomly generated circuits used by \sys, the evaluation
% pipeline can ingest circuits from established quantum benchmark suites:
% SupermarQ~\cite{tomesh2022supermarqscalablequantumbenchmark}, MQT
% Bench~\cite{Quetschlich_2023}, and QASM
% Bench~\cite{li2022qasmbenchlowlevelqasmbenchmark}. These suites provide
% externally curated circuits that complement the synthetic generator. The random
% generator is useful for broad, controlled coverage of the feature space, while
% benchmark-suite circuits provide recognizable algorithmic workloads for
% cross-checking that trends observed on generated circuits also appear on
% community benchmarks. After deduplicating equivalent algorithm families across
% suites and problem sizes, the benchmark-suite inventory covers 85 unique
% quantum algorithm families.

% InferQ includes standalone loaders under
% \path{scripts/benchmark_suites}. Each loader produces Qiskit
% \texttt{QuantumCircuit} objects and tags each circuit with a source-suite
% identifier and a benchmark name. The ingestion script
% \path{scripts/ingest_benchmarks.py} then passes these circuits through the
% same feature extraction, simulation, duplicate detection, hashing, and local
% storage path used for generated InferQ circuits. As a result, benchmark-suite
% circuits can be merged into the same metadata tables and evaluated using the
% same RDBMS, simulator-selection, and out-of-core workflows.

% \begin{table}[t]
%   \centering
%   \caption{External benchmark-suite support in the InferQ ingestion pipeline.}
%   \label{tab:benchmark-suite-ingestion}
%   \footnotesize
%   \begin{tabularx}{\columnwidth}{@{}lXX@{}}
%     \toprule
%     \textbf{Suite} & \textbf{Quantum Algorithms} & \textbf{Ingestion notes} \\
%     \midrule
%     SupermarQ & 8 benchmark entries & Converted from Cirq to Qiskit \\
%     MQT Bench & 34 benchmark entries & Loaded through the MQT Bench API \\
%     QASM Bench & 67 benchmark entries & Parsed from a pinned QASM snapshot \\
%     \bottomrule
%   \end{tabularx}
% \end{table}

% For SupermarQ, the loader covers the application benchmarks exposed by the
% suite, including GHZ, Hamiltonian simulation, Mermin-Bell, error-correction,
% QAOA, and VQE-style circuits. For MQT Bench, the loader uses a curated
% algorithm-level subset spanning state preparation, Fourier and phase-estimation
% circuits, oracle algorithms, variational circuits, quantum walks, random
% circuits, and representative arithmetic/error-correction circuits. For QASM
% Bench, InferQ vendors the small and medium tiers as OpenQASM files, records the
% upstream commit in a manifest, and excludes the large tier to keep the
% evaluation within the same qubit and runtime budgets as the rest of the paper.

% This design makes the external suites an additive evaluation source rather than
% a replacement for the generated corpus. Generated circuits remain necessary for
% systematic coverage, controllable scale, and reproducibility under fixed random
% seeds. The external suites provide an additional validity check by exercising
% InferQ on named circuits that are already used by the quantum-computing
% community.



% \begin{table*}[t]
%   \centering
%   \caption{Two-profile database-specific tuning settings used for the 7{,}705-circuit evaluation.}
%   \label{tab:finetuned-rdbms-profiles}
%   \footnotesize
%   \setlength{\tabcolsep}{4pt}
%   \begin{tabularx}{\textwidth}{@{}llXX@{}}
%     \toprule
%     \textbf{Engine} & \textbf{Profile} & \textbf{Parameters} & \textbf{Rationale} \\
%     \midrule
%     DuckDB
%       & \texttt{small\_size}
%       & \texttt{memory\_limit=10780MB}; \texttt{threads=2};
%         \texttt{preserve\_insertion\_order=false};
%         \texttt{max\_temp\_directory\_size=256GB}
%       & Gives DuckDB enough buffer-pool memory and limited parallelism for
%         short queries, while disabling insertion-order preservation to allow
%         spill-capable operators. \\
%     DuckDB
%       & \texttt{large\_size}
%       & \texttt{memory\_limit=7409MB}; \texttt{threads=1};
%         \texttt{preserve\_insertion\_order=true};
%         \texttt{max\_temp\_directory\_size=128GB}
%       & Reduces concurrent memory pressure and temp-space exposure for larger
%         query plans, where planner/runtime overhead can dominate. \\
%     SQLite
%       & \texttt{small\_size}
%       & \texttt{cache\_mb=64}; \texttt{db\_path=:memory:};
%         \texttt{mmap\_mb=0}; \texttt{temp\_store=MEMORY};
%         \texttt{threads=1}
%       & Keeps small workloads fully in memory and avoids filesystem overhead. \\
%     SQLite
%       & \texttt{large\_size}
%       & \texttt{cache\_mb=128}; file-backed database;
%         \texttt{mmap\_mb=128}; \texttt{temp\_store=MEMORY};
%         \texttt{threads=1}
%       & Uses a persistent database file and modest mmap window to stabilize
%         larger intermediate results while retaining memory-backed temporaries. \\
%     PostgreSQL
%       & \texttt{small\_size}
%       & \texttt{work\_mem=256MB}; \texttt{temp\_buffers=128MB};
%         \texttt{effective\_cache\_size=8GB};
%         \texttt{hash\_mem\_multiplier=2.0};
%         \texttt{max\_parallel\_workers\_per\_gather=4};
%         \texttt{join\_collapse\_limit=1}; \texttt{from\_collapse\_limit=1};
%         \texttt{jit=off}
%       & Allocates more per-query memory and parallelism to accelerate small
%         contractions, while preserving the emitted join order and avoiding JIT
%         overhead. \\
%     PostgreSQL
%       & \texttt{large\_size}
%       & \texttt{work\_mem=55MB}; \texttt{temp\_buffers=65MB};
%         \texttt{effective\_cache\_size=2GB};
%         \texttt{hash\_mem\_multiplier=1.94};
%         \texttt{max\_parallel\_workers\_per\_gather=0};
%         \texttt{join\_collapse\_limit=4}; \texttt{from\_collapse\_limit=1};
%         \texttt{jit=off}
%       & Caps per-node memory and disables parallel workers for large plans,
%         reducing worker amplification and making externalization behavior more
%         predictable. \\
%     Umbra
%       & comparison engine
%       & PostgreSQL-compatible server interface; reported version:
%         \texttt{PostgreSQL 18.1 compatible umbra v0.2-1665-gfeeb4bb5d}
%       & Broadens the engine comparison with a modern compiled analytical DBMS
%         that accepts the same SQL workload through a PostgreSQL-compatible
%         client path. \\
%     \bottomrule
%   \end{tabularx}
% \end{table*}
\begin{table*}[t]
  \centering
  \caption{Database-specific tuning settings used for the 7{,}705-circuit evaluation. 
  % The rationale for each profile is discussed in Section~\ref{app:mlos-runner}.
  }
  \label{tab:finetuned-rdbms-profiles}
  \footnotesize
  \setlength{\tabcolsep}{4pt}
  \begin{tabularx}{\textwidth}{@{}llX@{}}
    \toprule
    \textbf{Engine} & \textbf{Profile} & \textbf{Parameters} \\
    \midrule
    DuckDB
      & \texttt{small\_size}
      & \texttt{memory\_limit=10780MB}; \texttt{threads=2};
        \texttt{preserve\_insertion\_order=false};
        \texttt{max\_temp\_directory\_size=256GB} \\
    DuckDB
      & \texttt{large\_size}
      & \texttt{memory\_limit=7409MB}; \texttt{threads=1};
        \texttt{preserve\_insertion\_order=true};
        \texttt{max\_temp\_directory\_size=128GB} \\
    SQLite
      & \texttt{small\_size}
      & \texttt{cache\_mb=64}; \texttt{db\_path=:memory:};
        \texttt{mmap\_mb=0}; \texttt{temp\_store=MEMORY};
        \texttt{threads=1} \\
    SQLite
      & \texttt{large\_size}
      & \texttt{cache\_mb=128}; file-backed database;
        \texttt{mmap\_mb=128}; \texttt{temp\_store=MEMORY};
        \texttt{threads=1} \\
    PostgreSQL
      & \texttt{small\_size}
      & \texttt{work\_mem=256MB}; \texttt{temp\_buffers=128MB};
        \texttt{effective\_cache\_size=8GB};
        \texttt{hash\_mem\_multiplier=2.0};
        \texttt{max\_parallel\_workers\_per\_gather=4};
        \texttt{join\_collapse\_limit=1}; \texttt{from\_collapse\_limit=1};
        \texttt{jit=off} \\
    PostgreSQL
      & \texttt{large\_size}
      & \texttt{work\_mem=55MB}; \texttt{temp\_buffers=65MB};
        \texttt{effective\_cache\_size=2GB};
        \texttt{hash\_mem\_multiplier=1.94};
        \texttt{max\_parallel\_workers\_per\_gather=0};
        \texttt{join\_collapse\_limit=4}; \texttt{from\_collapse\_limit=1};
        \texttt{jit=off} \\
    % Umbra
    %   & comparison engine
    %   & PostgreSQL-compatible server interface; reported version:
    %     \texttt{PostgreSQL 18.1 compatible umbra v0.2-1665-gfeeb4bb5d} \\
    \bottomrule
  \end{tabularx}
\end{table*}


% \section{DBMS-Specific Tuning }
% \label{app:finetuned-rdbms}

% To understand how RDBMS configuration affects simulation performance, we have added  
% DBMS-specific tuning for PostgreSQL, DuckDB, and SQLite. 
% % We keep the input circuit set and \sys’s SQL generation
% % pipeline unchanged, and introduce a new runner that varies (i) \emph{DBMS configuration
% % parameters} (memory, parallelism, temporary I/O, and optimizer/planner settings) and
% % (ii) \emph{execution-structure choices} exposed by the SQL formulation (CTE
% % materialization and split execution). 
% The code and documentation are updated in the
% anonymous repository.\footnote{\url{}}.


% \medskip
% \para{Automated tuning via configuration search}
% We initially explored manual tuning to understand the impact of major parameters, but the
% configuration space is large, engine-specific, and not scalable across thousands of
% circuits. We therefore adopt an automated configuration-search workflow based on MLOS
% (Machine Learning for Operating Systems)\footnote{\url{https://github.com/microsoft/MLOS}}:
% the workload is fixed, and this tuner searches over engine parameters to minimize runtime.
% % The all-engine tuning experiment uses a 162-circuit validation subset, while the larger
% % SQLite-vs.-Qiskit comparison is run on the 7{,}705-circuit dataset used in
% % Figure~\ref{fig:aer-vs-sqlite-wins}. We keep these two roles separate: the 162-circuit
% % experiment selects and validates DBMS configurations, whereas the 7{,}705-circuit SQLite
% % run evaluates the main memory/runtime comparison at scale. For reproducibility, we
% % document the parameter spaces, seed profiles, and runner details in
% % Section~\ref{app:tuning-repro}.

% \medskip
% \para{Our workload-specific tuning}
% In preliminary experiments, we observed that a \emph{single} global configuration is suboptimal
% because circuits induce heterogeneous query shapes (e.g., very short vs.\ long CTE chains)
% and different bottlenecks (overhead-dominated vs.\ memory/IO-dominated). We therefore
% introduce a workload-specific separation. 
% for example 
% for the 162-circuit all-engine MLOS validation we divide circuits by median
% \texttt{circuit\_size} (33 gates): circuits with \texttt{circuit\_size} $\le 33$ use a
% \texttt{small\_size} profile (82 circuits), and the remaining circuits use a
% \texttt{large\_size} profile (80 circuits).
% This
% allows us to tune configurations to structurally different workloads: 
%  small circuits benefit from
%  settings like larger caches and limited parallelism, while larger
% circuits often require certain settings to reduce planner overhead and avoid
% unpredictable spilling.
% % For example, for the 7{,}705-circuit dataset\footnote{\url{https://anonymous.4open.science/r/InferQ-ADEF/analysis/training_data/rdbms_training_data.parquet}}, we use a separate
% % static-contraction SQLite profile: both size classes keep the database in memory
% % (\texttt{db\_path=:memory:}), use a 64~MB SQLite cache, disable cache spilling, and place
% % temporary structures in file-backed temp storage (\texttt{temp\_store=FILE}). 
% This run is
% not an additional MLOS search; it is the scaled SQLite validation used to compare against
% the best Qiskit Aer method. Table~\ref{tab:finetuned-rdbms-profiles} shows the dominating parameters found by MLOS after we set such a work-specific tuning. 
% % small circuits can benefit from more aggressive settings (e.g., caching), whereas
% % larger circuits often require choices to avoid planner overhead,
% % unfavorable materialization, and unstable spilling behavior.


% % In our SQL-based simulator, the tensor-contraction sequence is written as a chain of Common Table Expressions (CTEs), each encoding one contraction step (join + aggregate). Our preliminary finding is that data representation, SQL execution-plan overhead, and Aer method choice drive the observed differences.
% % We report CTE materialization because it changes how the optimizer/executor
% % treats the SQL query for each circuit (\texttt{WITH ... SELECT}) without converting the
% % workload into a sequence of explicit intermediate relations. In the out-of-core QFT
% % (3--26 qubits) experiment at a 4~GB cap, \texttt{monolithic\_materialized} preserves the
% % same completion range as \texttt{monolithic} for DuckDB (3--24 qubits) and for
% % SQLite/PostgreSQL (3--26 qubits), while reducing matched median runtime relative to
% % \texttt{monolithic} by roughly $2.1\times$ for DuckDB and $1.8\times$ for PostgreSQL;
% % SQLite is essentially unchanged. \texttt{split} remains the most robust out-of-core
% % structure, completing all PostgreSQL and SQLite configurations and all DuckDB
% % configurations except the 26-qubit/4~GB case.

% \medskip
 





% %https://github.com/microsoft/MLOS%

% \begin{table*}[t]
%   \centering
%   \caption{Fine-tuned all-engine validation on the 162-circuit MLOS subset. Each engine is run twice per circuit. Ratios compare the tuned engine to the best successful Qiskit Aer method; values below $1.0\times$ favor the RDBMS engine.}
%   \label{tab:finetuned-rdbms-results}
%   \small
%   \setlength{\tabcolsep}{4pt}
%   \begin{tabular}{lrrrrrrr}
%     \toprule
%     \textbf{Engine} & \textbf{Succ. circuits} & \textbf{Succ. runs} & \textbf{Timeout/error runs} & \textbf{Time wins} & \textbf{Median time} & \textbf{Memory wins} & \textbf{Median memory} \\
%     \midrule
%     SQLite & 136/162 & 270/324 & 54/324 & 7/136 & $4.99\times$ & 95/136 & $0.51\times$ \\
%     PostgreSQL & 134/162 & 267/324 & 57/324 & 0/134 & $12.00\times$ & 9/134 & $2.23\times$ \\
%     DuckDB & 71/162 & 139/324 & 185/324 & 0/71 & $33.68\times$ & 11/71 & $1.35\times$ \\
%     \bottomrule
%   \end{tabular}
% \end{table*}

% \begin{table}[t]
%   \centering
%   \caption{SQLite static-contraction validation on the 7{,}705-circuit dataset. Winner counts compare tuned SQLite, untuned baseline SQLite, and the best successful Qiskit Aer method.}
%   \label{tab:sqlite-7705-static-results}
%   \footnotesize
%   \setlength{\tabcolsep}{3pt}
%   \begin{tabular}{lrrr}
%     \toprule
%     \textbf{Objective} & \textbf{Tuned SQLite} & \textbf{Baseline SQLite} & \textbf{Qiskit Aer} \\
%     \midrule
%     Runtime & 0 (0.0\%) & 1{,}067 (13.8\%) & 6{,}638 (86.2\%) \\
%     Peak memory & 2{,}085 (27.1\%) & 1{,}818 (23.6\%) & 3{,}802 (49.3\%) \\
%     \bottomrule
%   \end{tabular}
% \end{table}

% \para{Results}
% Table~\ref{tab:finetuned-rdbms-results} shows that simulation results after tuning over the 162 circuit datasets\footnote{\url{https://anonymous.4open.science/r/InferQ-ADEF/analysis/training_data/rdbms_all_methods_training_data.parquet}}. We see that PostgreSQL and DuckDB's performance are significatnly improved by automatic tuning, esepcially for peak memeory, athough   Qiskit Aer remains the fastest backend for
% most general circuits.

% We did not observe SQLite's performance did not have an performance gain from autmoatic tuning, so we have conducted manual tuning. Even so, the runtime imporvement is mild 

% The results on 7{,}705-circuit results for SQLite similar. 
% Tuned SQLite improves over the untuned SQLite baseline in 4{,}946/7{,}404
% matched memory comparisons (66.8\%) and uses less peak memory than the best Qiskit Aer
% method in 3{,}891/7{,}404 matched comparisons (52.6\%). Considering the overall winner
% among tuned SQLite, baseline SQLite, and Qiskit Aer, SQLite is the peak-memory winner on
% 3{,}903/7{,}705 circuits (50.7\%): 2{,}085 wins come from tuned SQLite and 1{,}818 from
% the baseline SQLite configuration. Runtime shows the opposite pattern: Qiskit Aer remains
% the winner on 6{,}638/7{,}705 circuits (86.2\%), while all SQLite runtime wins come from
% the untuned baseline configuration (1{,}067/7{,}705 circuits). Thus, SQL-based simulation
% is not universally faster, but its sparse relational representation can be substantially
% more memory-efficient, making learned routing and DBMS-level physical optimization
% valuable.

% % The profile-level breakdown confirms that larger query plans are the harder regime. In
% % the \texttt{small\_size} profile (82 circuits; 164 runs per engine), SQLite completes
% % 153 runs and times out on 11, PostgreSQL completes 152 and times out on 12, and DuckDB
% % completes 97 and times out on 67. In the \texttt{large\_size} profile (80 circuits; 160
% % runs per engine), SQLite completes 117 runs and times out on 43, PostgreSQL completes
% % 115 runs with 44 timeouts and one error, and DuckDB completes 42 runs and times out on
% % 118. No configurations are skipped in this validation run; each engine is attempted for
% % each circuit under its selected profile. By objective, tuned RDBMSs are most helpful for
% % memory on both profiles: they beat Qiskit Aer on 61/78 completed \texttt{small\_size}
% % comparisons and 36/60 completed \texttt{large\_size} comparisons, but runtime wins occur
% % only in the smaller profile (7/78).

% % Table~\ref{tab:sqlite-7705-static-results} reports the actual SQLite static-contraction
% % results on the 7{,}705-circuit dataset. Tuned SQLite completes 7{,}404/7{,}705 circuits
% % (96.1\%) and times out on the remaining 301 circuits; at the run level, this corresponds
% % to 14{,}800 successful runs and 610 timeouts across two runs per circuit. The completion
% % rate is higher for the \texttt{small\_size} profile (3{,}829/3{,}889 circuits) than for the
% % \texttt{large\_size} profile (3{,}575/3{,}816 circuits), confirming that larger query plans
% % remain the harder regime.
\section{DBMS-Specific Tuning}
\label{app:finetuned-rdbms}

To understand how DBMS configuration affects simulation performance, we have added
DBMS-specific tuning for PostgreSQL, DuckDB, and SQLite.\footnote{We did not include Umbra in the DBMS-specific tuning study because its publicly available documentation for the released Docker image only describes basic access through the PostgreSQL-compatible server and command-line interface, and we could not identify a stable, documented set of engine-level tuning parameters needed to define a reproducible configuration-search space: \url{https://hub.docker.com/r/umbradb/umbra }} We keep the circuit set and
\sys’s SQL generation pipeline unchanged and only vary DBMS-level configuration.

\medskip
\para{Automated tuning via configuration search}
We initially explored manual tuning to understand the impact of major parameters, but the
configuration space is large, engine-specific, and not scalable across thousands of
circuits. We therefore adopt an automated configuration-search workflow based on MLOS
(Machine Learning for Operating Systems)\footnote{\url{https://github.com/microsoft/MLOS}}:
the workload is fixed, and the tuner searches over DBMS parameters to minimize runtime.
For reproducibility, we document the parameter spaces, seed profiles, and runner details in
Section~\ref{app:tuning-repro}.

\medskip
\para{Workload-specific tuning}
In preliminary experiments, we observed that a single global configuration is suboptimal:
circuits induce heterogeneous query shapes (e.g., short vs.\ long CTE chains) and
different bottlenecks (overhead-dominated vs.\ memory- and I/O-dominated). We therefore
introduce workload stratification by circuit size. For the 162-circuit all-engine MLOS
validation set,\footnote{\url{https://github.com/InfiniData-Lab/InferQ/blob/main/analysis/training_data/rdbms_all_methods_training_data.parquet}}
we split circuits by the median \texttt{circuit\_size} (33 gates): circuits with
\texttt{circuit\_size} $\le 33$ use a \texttt{small\_size} profile (82 circuits) and the
remaining circuits use a \texttt{large\_size} profile (80 circuits). This allows the tuner
to select configurations appropriate for structurally different workloads of small and large circuits.
% benefit from settings such as larger caches and limited parallelism, whereas larger
% circuits often require more conservative settings to reduce planner overhead and avoid
% unpredictable spilling. 
Table~\ref{tab:finetuned-rdbms-profiles} summarizes the final
profiles. We observed that SQLite did not obtain much performance gain from automatic tuning, so we have tuned it manually\footnote{Tuning results for the 7705 dataset can be found here: \url{https://github.com/InfiniData-Lab/InferQ/blob/main/analysis/finetuned/rdbms7705_sqlite_static_contr/rdbms7705_sqlite_static_contr_results.csv}}.

\begin{table}[t]
  \centering
  \caption{Tuned DB performance on the 162 circuits. Time and peak-memory wins are counted against the best successful Qiskit Aer method per circuit.}
  \label{tab:finetuned-rdbms-results}
  \small
  \setlength{\tabcolsep}{6pt}
  \begin{tabular}{lrr}
    \toprule
    \textbf{Backend} & \textbf{Time wins} & \textbf{Memory wins} \\
    \midrule
    SQLite      & 34/162  & 108/162 \\
    DuckDB      & 0/162   & 0/162 \\
    PostgreSQL  & 0/162   & 0/162 \\
    Umbra       & 0/162   & 0/162 \\
    \midrule
    RDBMS total     & 34/162  & 108/162 \\
    Qiskit Aer total & 128/162 & 54/162 \\
    \bottomrule
  \end{tabular}
\end{table}

\para{Results}
Table~\ref{tab:finetuned-rdbms-results} reports results after tuning 
% \footnote{Tuni \url{https://anonymous.4open.science/r/InferQ-ADEF/analysis/finetuned/rdbms162_all_engines_120s/rdbms162_all_engines_120s_results.csv}}
on the 162-circuit dataset, which are also the results for Figure~\ref{fig:rdbmsqiskitcomparison}. Overall, Qiskit Aer
remains the runtime winner on most circuits (128/162). Among the evaluated DBMS backends,
only SQLite achieves runtime wins (34/162); DuckDB, PostgreSQL, and Umbra do not win on
runtime on this subset. In contrast, SQLite is the peak-memory winner on 108/162 circuits ($66.7\%$),
while Qiskit Aer wins peak memory on 54/162 circuits.
We have also evaluated the tuned configurations on the full \num{7705}-circuit dataset
(Figure~\ref{fig:aer-vs-sqlite-wins}). The result is similar: RDBMS backends achieve the
lowest runtime on \num{1101} circuits (14.3\%), and all of these runtime wins come from the SQLite, for which tuning did not have much improvement. For peak memory, RDBMS backends win on \num{3902}
circuits (50.6\%), dominated by tuned SQLite (\num{3901} wins) with one additional win
by Umbra. Relative to the untuned setting, tuning increases SQLite's peak-memory win count
from \num{3196} to \num{3901} circuits.

\medskip
  \vspace{0.2cm}
\noindent
\setlength{\fboxsep}{6pt}
\fcolorbox{black!15}{black!4}{%
  \parbox{\dimexpr\linewidth-2\fboxsep-2\fboxrule\relax}{
\para{Take-away}
Taken together with Appendix~A (out-of-core) and Appendix~F (memory profiling), these
results suggest two concrete future optimization opportunities for DBMS-backed simulation. 
\\i) For
\emph{large} circuits where memory is the bottleneck, DBMS execution is uniquely useful
because it can externalize operator workspaces and intermediate relations. For such simulation,
end-to-end performance is largely driven by spill and intermediate growth, motivating
spill-aware planning (join ordering, operator selection, and materialization control) and  physical design studies over
intermediate tables (indexes, partitioning and layout, and caching reused intermediates) as discussed in Appendix~H.
\\ii) For \emph{small} circuits where Aer finishes in milliseconds, DBMS runtime is often
dominated by fixed overheads (SQL compilation, planning, and repeated query execution),
motivating workload-level optimizations such as batching,  reuse, and
handling repeated subcircuit computations across many runs. 
\\
\sys is a convenient tool
for future database research in these directions, because it generates controlled workloads
(small or large, sparse or dense) and emits database-native SQL along with query-shape and
circuit metadata needed for systematic evaluation. 
}}
%   \begin{table}[t]
%     \centering
%     \caption{Tuned DB performance on the 162 circuits. Time and
%   memory wins are compared to the best successful Qiskit Aer method.}
%     \label{tab:finetuned-rdbms-results}
%     \small
%     \setlength{\tabcolsep}{6pt}
%     \begin{tabular}{lrr}
%       \toprule
%       \textbf{Engine} & \textbf{Time wins} & \textbf{Memory wins} \\
%       \midrule
%       SQLite & 34/162 & 106/162 \\
%       PostgreSQL & 0/162 & 11/162 \\
%       DuckDB & 0/162 & 13/162 \\
%       \bottomrule
%     \end{tabular}
%   \end{table}



% \para{Results}
% Table~\ref{tab:finetuned-rdbms-results} reports results\footnote{Tuning results for the 162-circuit dataset can be found here: \url{https://anonymous.4open.science/r/InferQ-ADEF/analysis/finetuned/rdbms162_all_engines_120s/rdbms162_all_engines_120s_results.csv}}
% for the tuned configurations on the 162-circuit validation set. Overall, Qiskit Aer
% remains the runtime winner on most circuits. The performance of PostgreSQL and DuckDB has
% improved substantially, but they do not achieve any runtime wins against the best
% successful Aer method on this subset. SQLite runtime does not improve significantly
% either. We note that for many of these circuits, Qiskit Aer completes each simulation in
% milliseconds, so further gains may require optimizing the end-to-end simulation workload rather than tuning each query separately.

% The primary benefit of tuning is improved memory efficiency: tuned SQLite wins peak memory
% on 106/162 circuits (up from 94 before tuning), and PostgreSQL and DuckDB improve from 0
% to 11/162 and 13/162 (both 0 before tuning) peak-memory wins, respectively.

% Over the large dataset of \num{7705} circuits, the results are similar. SQLite remains to be over fastest engine among RDBMSs. Our tuning did ot change much of its perfercentage of winning in runtime. But for memory, before tuning SQLite wins 3,196 circuits (41.6\%), after tuning 3,901. Umbra also wins 1 circuit. In totoal, RDBMS wins 3,902 50.6\%, around half of case. This makes RDBMS a promising simulation tool in terms of memory.

% \para{Take-away} Query optimization is foundamental but hard problem. I believe \sys shows quantum circuit simulation can be a new type of SQL query workload to optimze
% From these resutls together, we can see there is much things to explore, especially runtime. Combing the results from Appendix~\ref{app:qft-cap-sensitivity}, we think it is intersting to explore 1) optimize queries for single large circuit that requires performance degrade gracefully out of core; and 2) in quantum experiments they often run multiple simulation or same simulation repeatly. Then how to optimze a workload of multiple SQL queries representing such a simulation workload. \sys is a convient tool for generating SQL queries for such future research.
     


% \begin{table}[t]
%   \centering
%   \caption{SQLite static-contraction validation on the 7{,}705-circuit dataset. Winner counts compare tuned SQLite, untuned baseline SQLite, and the best successful Qiskit Aer method.}
%   \label{tab:sqlite-7705-static-results}
%   \footnotesize
%   \setlength{\tabcolsep}{3pt}
%   \begin{tabular}{lrrr}
%     \toprule
%     \textbf{Objective} & \textbf{Tuned SQLite} & \textbf{Baseline SQLite} & \textbf{Qiskit Aer} \\
%     \midrule
%     Runtime & 0 (0.0\%) & 1{,}067 (13.8\%) & 6{,}638 (86.2\%) \\
%     Peak memory & 2{,}085 (27.1\%) & 1{,}818 (23.6\%) & 3{,}802 (49.3\%) \\
%     \bottomrule
%   \end{tabular}
% \end{table}

% For the larger 7{,}705-circuit evaluation, Table~\ref{tab:sqlite-7705-static-results}
% shows the same pattern at scale.\footnote{\url{https://anonymous.4open.science/r/InferQ-ADEF/analysis/training_data/rdbms_training_data.parquet}}
% Qiskit Aer remains the runtime winner on 6{,}638/7{,}705 circuits (86.2\%). In contrast,
% SQLite is competitive on peak memory: tuned SQLite is the peak-memory winner on
% 2{,}085/7{,}705 circuits, and the baseline SQLite configuration wins on 1{,}818/7{,}705
% circuits. These results confirm that SQL-based simulation is not universally faster, but
% its sparse relational representation can be substantially more memory-efficient, which
% motivates learned routing and DBMS-level optimization for circuit-dependent performance.



\section{Indexes, partitioning, and CTE query structure}
\medskip
 
In \sys, each circuit is compiled into a \emph{single SQL statement} of the form
\texttt{WITH ... SELECT}, where each tensor contraction step is represented as a Common Table Expressions (CTE) and
implemented as a join followed by an aggregate, consistent with prior SQL-based
formulations \cite{hai2025quantum, einsteinSQL2023}. This CTE-chain representation is
useful because it makes the contraction order explicit and exposes the workload to the
DBMS optimizer as a sequence of standard relational operators (joins and group-bys).
However, it also limits classic physical design actions: intermediate CTE results are
 not persistent relations, so users cannot directly create secondary indexes, define
table partitioning, or maintain materialized views over those intermediates. Instead, the
engine may build transient internal structures (e.g., SQLite's automatic indexes and temporary B-tree structures), but these are optimizer-managed and not user-controlled. To enable explicit
physical design, we added a second query generator which lowers the same contraction pipeline
into a sequence of \texttt{CREATE TEMP TABLE AS} statements; each intermediate becomes a
first-class relation on which indexes and partitioning can be applied. 
Note that this is not part of \sys's standard pipeline; we added it to support the
experiments in this section. Systematically exploring index and partition selection for
these intermediate relations is a natural direction for future work.
%

 

\medskip
\para{Implications for physical design}
The tuned profiles in Table~\ref{tab:finetuned-rdbms-profiles} provide a reproducible
baseline that isolates DBMS-level execution effects (memory budgeting, temporary-workspace
management, parallel execution, and planner constraints) for this workload. Combined with
\texttt{split} execution, they also define a concrete path for future physical design
studies using standard database methodology. For example, one can treat intermediate
tables as design objects and evaluate: (i) index selection on join keys for contraction
steps (e.g., B-tree or hash indexes depending on the engine and operator), (ii)
partitioning and layout decisions to improve join locality and reduce spill (where the engine
supports partitioned tables and partitionwise planning), and (iii) selective
materialization and caching of frequently reused intermediates (e.g., repeated subcircuits)
as materialized relations. Because \sys exposes circuit features and query-shape metadata
(e.g., number of contraction steps (CTEs), runtime, and spill proxies), these choices can be
studied systematically as physical design problems across engines and workload regimes.

\medskip
\para{Preliminary Results}
We have conducted an experiment similar to Appendix~\ref{app:qft-cap-sensitivity} under a
4~GB memory limit. We choose a dense circuit from existing benchmarks, namely QFT, and vary
its qubit count from 3 to 26. Except for DuckDB at 26 qubits, all simulations completed.
We observe that the new query generator, which forces CTE materialization, reduces runtime
by roughly $2.1\times$ for DuckDB and $1.8\times$ for PostgreSQL; SQLite is essentially
unchanged. This shows that CTE materialization can be a potentially effective approach
for improving DBMS-backed simulation on dense workloads under tight memory limits.

% \medskip
% \para{Summary}
% Overall, this DBMS-specific track complements our untuned evaluation by (i) quantifying
% how sensitive SQL-based simulation is to engine configuration, and (ii) separating
% single-statement CTE execution from split execution, which is the appropriate regime for
% explicit indexes, partitioning, and materialized intermediates. The resulting tuned
% profiles and runners provide a reproducible baseline for both performance evaluation and
% future physical-design research on database-native circuit simulation.

% \para{Indexes, partition, and CTEs} In this work, for each circuit generated by \sys,  
% we compile it into a  \texttt{WITH ... SELECT} statement whose contraction
% steps are represented by a single CTE chain in the SQL query following existing work \cite{hai2025quantum, einsteinSQL2023}. 
% % Now we see that this query structure limits the opportunities for 
% % user-managed physical design: there are no durable intermediate relations on which one
% % can create explicit indexes or partitions. Instead, it relies on each
% % engine's internal choices, e.g., SQLite's automatic indexes and temporary B-tree
% % structures. In contrast, \texttt{split} execution lowers the same contraction plan into
% % a sequence of \texttt{CREATE TEMP TABLE AS} statements, making each intermediate an
% % explicit physical relation. This is the natural regime to study user-managed indexes,
% % materialized intermediates, and (where supported) partitioning planning.
% This matters for physical design. In this query
% shape there are no durable intermediate relations on which \sys can build
% user-managed indexes. Explicit indexes are therefore only applicable to split
% query mode, where the same CTE chain is lowered into per-step
% \texttt{CREATE TEMP TABLE} statements and each intermediate can be treated as a
% physical relation. In the monolithic profile we instead rely on each engine's
% internal planning choices, including automatic or ephemeral index structures
% where available, such as SQLite's automatic indexes. Partitioning has the same
% scope: PostgreSQL partitionwise planner switches can be enabled, but the
% monolithic CTEs are not declared partitioned tables. Partitioning is therefore a
% future split-mode or materialized-table design choice rather than a central knob
% in the two-profile monolithic experiment.

% The fine-tuned runner uses the raw monolithic IQS query shape: each circuit is
% compiled into a single \texttt{WITH ... SELECT} statement whose contraction
% steps are represented as CTEs. This matters for physical design. In this query
% shape there are no durable intermediate relations on which InferQ can build
% user-managed indexes. Explicit indexes are therefore only applicable to split
% query mode, where the same CTE chain is lowered into per-step
% \texttt{CREATE TEMP TABLE} statements and each intermediate can be treated as a
% physical relation. In the monolithic profile we instead rely on each engine's
% internal planning choices, including automatic or ephemeral index structures
% where available, such as SQLite's automatic indexes. Partitioning has the same
% scope: PostgreSQL partitionwise planner switches can be enabled, but the
% monolithic CTEs are not declared partitioned tables. Partitioning is therefore a
% future split-mode or materialized-table design choice rather than a central knob
% in the two-profile monolithic experiment.


 

\section{Additional Configuration Details on Experiments for Reproducibility}
\label{app:tuning-repro}

This section documents how we measure spilling performance and how we execute tuned runs.

\subsection{Circuit Construction Details}
\label{app:expander-details}
We design a circuit family with large qubit
counts but fixed sparse values.
Each circuit is parameterized by $(n,h)$, where $n$
is the number of qubits and $h<n$ is a small seed size. The construction applies Hadamards
only to the $h$ seed qubits and then applies only CNOT layers, which preserve support
size; therefore the final state has exactly $2^h$ non-zero amplitudes in the computational
basis. We
evaluate a sample of Expander circuits with $n \in [40,45]$, where a dense statevector
would require approximately 16--512~TB.  
An Expander circuit on $n$ qubits is parameterized by a seed size $h<n$. We apply
Hadamards to the first $h$ qubits to create a superposition over $2^h$ basis states and
initialize the remaining qubits to $\ket{0}$. We then apply only CNOT layers, which are
reversible linear transformations and preserve the number of basis states in support.
Consequently, the final state remains a uniform superposition over an affine subspace of
dimension $h$ and has exactly $2^h$ non-zero amplitudes in the computational basis (each
with magnitude $2^{-h/2}$). Figure~\ref{fig:expander_circuit} illustrates the circuit
structure.

\begin{figure}[t]
    \centering
    \includegraphics[width=1.0\linewidth]{images/expander_circuit_diagram.pdf}
    \caption{Expander circuit structure: Hadamards on $h$ seed qubits followed by CNOT-only
    layers that expand connectivity while preserving sparse support.}
    \label{fig:expander_circuit}
\end{figure}

Such sparse or highly structured states occur in practice in affine and stabilizer-style subroutines and
in workloads where computation is restricted to a low-dimensional subspace. Moreover,
even when the final state is dense, intermediate tensors and relations during contraction
can be much smaller than $2^n$. Separately, amplitude-amplification procedures (e.g.,
Grover-style iterations) tend to \emph{concentrate probability mass} onto a small subset
of basis states as they iterate, and sparse Hamiltonians and sparse linear systems (as in
eigensolver and HHL-style routines) often induce structured tensors. These regimes
motivate studying DBMS backends with sparse representations and out-of-core execution.

\subsection{Spill Configuration for Out-of-core experiments}
\label{app:spill-instrumentation}

We run SQL queries on PostgreSQL, DuckDB, and SQLite, sweeping Docker memory limit, e.g.,  16~GB. 
All containers use a one-CPU quota, swap is disabled by setting Docker
\texttt{--memory-swap} equal to the memory cap, and host page cache is not
explicitly dropped between runs. DuckDB uses one thread and sets
\texttt{memory\_limit} to the cap minus a 1024~MB headroom pad (3072, 7168,
and 15360~MB for the 4, 8, and 16~GB caps, respectively). SQLite uses a
64~MB cache budget for both main and temporary schemas. PostgreSQL uses the
tuned PostgreSQL~12.22 container, with the server container capped at the
experiment memory limit and \texttt{temp\_file\_limit} set to 64~GB. Spill is
measured as \texttt{postgres\_explain\_temp\_written} for PostgreSQL and as
\texttt{proc\_io\_write\_bytes} for DuckDB and SQLite.

We report spill volume using engine-specific sources when available, and otherwise use a
conservative operating system level proxy.

\noindent\textbf{PostgreSQL.} We compute exact temporary spill bytes from \texttt{EXPLAIN} (ANALYZE, BUFFERS, FORMAT JSON) by converting the reported temp written/read blocks to bytes using the 8192-byte block size; we treat this value as the per-query externalization metric.

\noindent\textbf{DuckDB.} We use \texttt{proc\_\allowbreak io\_\allowbreak write\_\allowbreak bytes} as the primary spill/workspace proxy, and when available we additionally parse DuckDB profiling JSON fields such as \texttt{temporary\_\allowbreak storage\_\allowbreak bytes} and \texttt{spilled\_\allowbreak bytes} for diagnostics only, since the profile schema can vary across versions.

\noindent\textbf{SQLite.} Similar to DuckDB, here we use \texttt{proc\_\allowbreak io\_\allowbreak write\_\allowbreak bytes} as the spill/workspace proxy. SQLite does not expose a stable query-level temp-byte counter through Python’s \texttt{sqlite3} interface, so we intentionally do not report a DB-native temp metric.

% \noindent\textbf{Qiskit Aer.} Aer does not spill, so we report no spill metric and instead record the simulator method, wall time, status, cgroup peak RSS, and whether execution fails due to memory.

\subsection{Configuration for DB tuning}
\label{app:mlos-runner}
We add a dedicated runner under \texttt{scripts/finetuned\_rdbms} that reuses the same
InferQ circuits and SQL generation path, but executes the emitted SQL under
engine-specific configurations rather than untuned defaults. The same directory contains
an automated tuning script (\texttt{mlos\_tune\_rdbms.py}) based on MLOS.

MLOS keeps the workload fixed (each circuit is compiled once into a SQL
query) and searches only over DBMS configuration parameters. For each engine, the tuner
evaluates seed profiles (e.g., \texttt{balanced}, \texttt{fast}, \texttt{spill\_safe})
and then uses a black-box optimizer backend (e.g., FLAML or SMAC) to propose additional
configurations. Its objective is the median runtime on a pilot set; timeouts and
execution failures receive a large penalty. All trials are recorded in
\texttt{mlos\_trials.csv}, and the best configuration per engine is stored in
\texttt{mlos\_best\_configs.json}. These winning configurations are then validated on the
full discovered circuit set via the fine-tuned runner using \texttt{--tuning-json}. Thus,
the reported improvements correspond to \emph{DBMS-level physical execution choices}
rather than changes to the circuit generator or SQL compilation.




Operationally, the runner generates SQL once per circuit, records query-shape metadata
(e.g., number of CTEs/contraction steps), and then executes the workload under the
selected engine profiles with warm-up and timed runs, configurable timeouts, and resume
support. To reduce total tuning cost, we optionally use SQLite as a lightweight gate:
if SQLite times out for a circuit under the tuned profile, we mark other engines as
\texttt{skipped\_sqlite\_timeout} for that circuit (i.e., we do not spend additional
compute on circuits already rejected by the gate). This produces a controlled comparison
of tuned physical execution while keeping the workload and SQL generation fixed.

% \para{Rationale for the two-profile settings.}
Table~\ref{tab:finetuned-rdbms-profiles} reports the best configurations selected and validated by our tuning workflow. For DuckDB, \texttt{small\_size} allocates a larger memory budget and modest parallelism for short query plans, while disabling insertion-order preservation to enable spill-capable operator choices; \texttt{large\_size} reduces memory and parallelism to limit concurrent memory pressure and temp-space exposure for deeper plans where planner/executor overhead can dominate. For SQLite, \texttt{small\_size} keeps the database and temporaries fully in-memory to avoid filesystem overhead, whereas \texttt{large\_size} uses a file-backed database with a modest mmap window to stabilize larger intermediate results while still keeping temporaries in memory. For PostgreSQL, \texttt{small\_size} increases per-node memory (\texttt{work\_mem}, \texttt{temp\_buffers}) and enables parallel workers to accelerate small contractions, while constraining join enumeration (\texttt{join\_collapse\_limit}/\texttt{from\_collapse\_limit}) to preserve the emitted join order and disabling JIT; \texttt{large\_size} caps per-node memory and disables parallel gather to reduce worker amplification and make externalization behavior more predictable. Finally, we include Umbra via its PostgreSQL-compatible frontend to broaden the engine comparison to a modern compiled analytical DBMS without changing the SQL workload.

\section{Generated Circuit Characterization}
\label{app:circuit-characterization}

\begin{figure*}[t]
    \centering
    \begin{subfigure}[t]{0.48\textwidth}
        \centering
        \includegraphics[width=\linewidth]{images/01_gate_mix.png}
        \caption{Top-20 gate-type distribution.}
        \label{fig:gate-mix}
    \end{subfigure}\hfill
    \begin{subfigure}[t]{0.48\textwidth}
        \centering
        \includegraphics[width=\linewidth]{images/03_depth_vs_width.png}
        \caption{Depth/width coverage.}
        \label{fig:depth-width}
    \end{subfigure}
    \caption{Generated-corpus gate mix and depth/width coverage. Color in \subref{fig:depth-width} indicates the number of circuits in each bin on a logarithmic scale.}
    \label{fig:generated-characterization-overview}
\end{figure*}

\begin{figure*}[t]
    \centering
    \includegraphics[width=\textwidth]{images/02_circuit_structure.png}
    \caption{Generated-corpus static circuit-feature distributions. Dashed lines mark medians.}
    \label{fig:circuit-structure}
\end{figure*}

\begin{figure*}[t]
    \centering
    \includegraphics[width=\textwidth]{images/04_graph_features.png}
    \caption{Generated-corpus interaction-graph feature distributions. Dashed lines mark medians.}
    \label{fig:interaction-graph}
\end{figure*}


To make the generated workload distribution transparent and reproducible, \sys records per-circuit metadata for each generated circuit and aggregates these fields over the full corpus. Figures~\ref{fig:generated-characterization-overview}, \ref{fig:circuit-structure}, and~\ref{fig:interaction-graph} summarize the resulting 202{,}975-circuit inventory from complementary perspectives: gate mix, static circuit structure, width/depth coverage, and interaction-graph structure. Together, these summaries show that the generator produces a broad workload rather than a collection of circuits that differ only in size.

\medskip
\para{Gate-mix distribution}
Figure~\ref{fig:gate-mix} reports the top-20 gate types by share of total gate count. The corpus has a substantial entangling component: \texttt{cx} alone accounts for 26.6\% of all gates, alongside common one-qubit rotations and Clifford-style operations such as \texttt{ry}, \texttt{x}, \texttt{h}, and \texttt{rz}. The remaining top gates include controlled-phase, controlled-Z, multi-controlled-X, swap, and higher-controlled variants, indicating that the generated workload includes both simple local transformations and larger multi-qubit control patterns.


\medskip
\para{Static circuit structure}
Figure~\ref{fig:circuit-structure} expands the view from gate identities to circuit-level structure. The generated corpus spans small and moderate-width circuits, with median values of 14 qubits, width 16, depth 59, and circuit size 190. It also varies substantially in two-qubit-gate count and two-qubit-gate percentage, which are important because they influence both tensor contraction structure and the amount of intermediate state growth during simulation. Additional static descriptors such as Pauli-gate count, locality ratio, idling score, and density score capture whether a circuit is mostly local or contains broader qubit interactions and idle regions. These distributions make the benchmark more informative than a width/depth grid alone.

\medskip
\para{Depth/width coverage}
Figure~\ref{fig:depth-width} shows the joint coverage over number of qubits and circuit depth. The heat map reports the number of circuits in each \emph{(number of qubits, depth)} bin using a logarithmic color scale, exposing both dense regions of the generated space and sparsely populated edge cases. The workload includes shallow circuits, deep circuits with depths above $10^3$, and circuit sizes extending to roughly 30 qubits in this corpus. This joint view is useful for interpreting performance results because simulator and DBMS behavior can change sharply when both width and depth increase, even when either feature alone appears moderate.

\medskip
\para{Interaction-graph structure}
Figure~\ref{fig:interaction-graph} characterizes the qubit interaction graphs induced by the generated circuits. We construct an interaction graph by treating qubits as vertices and adding an edge when two qubits participate in a multi-qubit operation. The resulting features include edge count, maximum and average degree, clustering coefficient, shortest-path statistics, diameter, radius, cut-related measures, central-point dominance, and adjacency-matrix variation. These graph features capture structure that is not visible from gate count, width, or depth alone: two circuits with the same number of qubits and a similar depth can still differ substantially in interaction density, path length, graph diameter, or degree distribution. Such differences affect SQL query shape, contraction opportunities, intermediate-result sizes, and simulator behavior.

Overall, these summaries make the generated workload transparent, reproducible, and analytically useful. With \sys, users can inspect the concrete properties that shape simulation performance, including gate composition, static circuit size, depth/width coverage, and qubit-interaction topology. This allows engine comparisons to be interpreted in terms of the circuits being simulated and helps identify which circuit families require particular simulation strategies.

 
\section{Subcircuit generation}
\label{sssec:circuittemplate}
In Section~\ref{sec:benchmark} we mentioned how to generate a subcircuit. Here, we provide more details.
% In Section~\ref{sec:benchmark} we mentioned that for each template $f$, we have parameter domain $\Theta_f$. Here we explain the below parameters that included in $\Theta_f$. From $\Theta_f$, \sys decides the value of the paramer through sampling. For instance, suppose a user sets \texttt{min\_depth}=5 and \texttt{max\_depth}=10.
% These bounds apply to each generated \emph{subcircuit} independently.
% When a template $f$ is selected, \sys forms the domain of its depth parameter $d$
% by intersecting the template’s intrinsic depth range with $[5,10]$.
% \sys then samples a value $d$ from this per-subcircuit domain and instantiates
% a subcircuit of exactly that depth.
% The total depth of the final circuit is the sum of the depths of all instantiated
% subcircuits.

% \[
% (n_q, d, k, r, \xi),
% \]
% where:
% \begin{itemize}
%     \item $n_q$ is the number of qubits;
%     \item $d$ is a target depth;
%     \item $k$ is the number of qubits used for evaluation tasks;
%     \item $r$ is a repetition count for repeatable subroutines inside a generator;
%     \item $\xi$ denotes a collection of \emph{template-specific configuration
%           variables}.
% \end{itemize}

% The variable $\xi$ captures auxiliary design choices that are local to a given
% circuit generator and are not shared across the global benchmark configuration.
% These include discrete architectural options, boolean feature flags, local layer
% counts, and generator-specific control parameters. All such variables are
% sampled by family-specific parameter generators and are constrained only by the
% global hyperparameters and the current circuit context.
% %
% Typical examples of variables contained in $\xi$ include whether a Quantum
% Fourier Transform includes final qubit-reversal swaps, the choice of feature map
% and ansatz type in variational circuits, the number of layers in QAOA-like
% subroutines, or the number of steps and coin operators in quantum walk circuits.

% Once a template is selected, a structural circuit template is first instantiated using the current instance configuration and the generator’s internal settings. This defines the qubit layout, gate topology, repetition structure, and depth contribution of the subcircuit. If the generator includes parametric gates, numerical gate values are then sampled from generator-specific distributions. This staged construction separates architectural variability from numerical variability, allowing \sys to explore both circuit structure and parameter space in a controlled and independent manner.




Each selected circuit template \(f\) is instantiated by sampling parameters from its
parameter domain \(\Theta_f\). In \sys, we write one sampled
parameter tuple as
\begin{equation}
\theta = (n_q, d, k, r, \xi) \in \Theta_f,
\end{equation}
where \(n_q\) is the qubit count, \(d\) is the depth  allocated to this subcircuit,
\(k\) is the evaluation-qubit, \(r\) is the repetition count for repeatable template
sections, and \(\xi\) is a tuple of template-specific configuration variables.

The domain \(\Theta_f\) is derived from the  parameters  in Table~\ref{tab:base_params}. For example, the per-subcircuit depth \(d\) is sampled uniformly from the defined range between  min\_depth and max\_depth.

The template-specific
configuration variables \(\xi\) captures local design choices supported by \(f\).
Typical examples include whether \texttt{QFT} includes final swaps, the choice of feature map
and ansatz in variational templates, the number of layers in QAOA-style templates, or the
number of steps and coin operators in quantum-walk templates.



\paragraph{Structural vs. numerical parameters.}
In \sys, template instantiation is staged. \sys first fixes the \emph{structure} of the subcircuit
(e.g., qubit assignment, gate topology, repetition pattern, and depth contribution) from
\((n_q,d,k,r,\xi)\). If the template contains parametrized gates,  \sys then samples the
\emph{numerical} values (e.g., rotation angles) from template-specific distributions. This
separation lets \sys explore structural diversity independently from numerical variability.



% \section{Attempt}
% Let $\mathcal S$ denote the space of subcircuits, and let $F$ be a finite set of circuit templates (Table~\ref{tab:base_params}).
% Let $B$ denote the \emph{context space}, i.e., the set of generation states maintained by \sys
% (e.g., remaining qubit/depth budgets, allocated registers, and composition constraints; Table~\ref{tbl:comb}).
% A \emph{circuit template} $f\in F$ is a parameterized generator
% \[
% f:\Theta_f \times B \to \mathcal S.
% \]
% At step $t$, after selecting a template $f_t$ and parameters $\theta_t\in\Theta_{f_t}$, \sys instantiates the corresponding
% subcircuit as
% \[
% S_t = f_t(\theta_t, b_{t-1}),
% \]
% where $b_{t-1}\in B$ is the current context.

% A circuit instance is an ordered composition of subcircuits:
% \[
% C \;=\; S_1 \circ S_2 \circ \cdots \circ S_T,
% \]
% where $T$ is determined by a stopping rule and/or budget exhaustion.

% \sys generates a template sequence $f_1,\dots,f_T$ incrementally.
% At step $t$, it samples the next template from a context-dependent categorical distribution on $F$:
% \[
% p^{(t)}(f)
% =
% \frac{\pi(f)\, w_t(f)\, \mathsf{feas}(b_{t-1}, f)}
% {\sum_{g\in F}\pi(g)\, w_t(g)\, \mathsf{feas}(b_{t-1}, g)},
% \qquad f\in F,
% \]
% where $\pi$ is a base distribution over templates,
% $w_t$ is a synergy weight (defined below),
% and $\mathsf{feas}(b,f)\in\{0,1\}$ is a hard feasibility predicate given context $b$
% (budgets and structural constraints).

% After sampling $f_t\sim p^{(t)}$, \sys samples $\theta_t\in\Theta_{f_t}$, instantiates
% $S_t=f_t(\theta_t,b_{t-1})$, appends $S_t$ to the circuit, and updates the context:
% \[
% b_t = \mathsf{Upd}(b_{t-1}, f_t, \theta_t, S_t).
% \]
% Here $\mathsf{Upd}$ deterministically updates remaining budgets (e.g., decreases remaining depth),
% tracks allocated registers, and records any state needed by later feasibility checks.

% To bias generation toward semantically coherent compositions (Table~\ref{tbl:comb}),
% \sys uses a set $\mathcal R$ of synergy rules.
% A rule is a triple $(T,U,\beta)$ with trigger set $T\subseteq F$,
% target set $U\subseteq F$, and multiplicative factor $\beta>0$.
% After selecting $f_t$, we update the synergy weights for the next step by
% \[
% w_{t+1}(f)
% =
% \prod_{(T,U,\beta)\in\mathcal R:\; f_t\in T,\; f\in U}\beta,
% \qquad f\in F,
% \]
% and then sample $f_{t+1}$ using $p^{(t+1)}$, i.e., after renormalizing over feasible templates.
% This keeps the process stochastic (diverse circuits) while increasing the probability of coherent transitions.

% The induced stochastic process on $(b_t,f_t)$ is first-order Markov: $(b_{t+1},f_{t+1})$ depends on the past only through $(b_t,f_t)$.


% \section{Non-Interleaved}
% At step $t$, \sys maintains a categorical distribution $p_t$ over templates.
% Given the current context $b_{t-1}$, it proceeds in two steps:

% \smallskip
% \noindent\textbf{(1) Sample a template.}
% First apply feasibility masking and renormalize:
% \[
% \bar p_t(f)
% =
% \frac{p_t(f)\,\mathsf{feas}(b_{t-1},f)}
% {\sum_{g\in F} p_t(g)\,\mathsf{feas}(b_{t-1},g)},
% \qquad f\in F,
% \]
% then sample $f_t\sim \bar p_t$.
% Next sample $\theta_t\in\Theta_{f_t}$, instantiate $S_t=f_t(\theta_t,b_{t-1})$,
% append $S_t$, and update the context:
% \[
% b_t=\mathsf{Upd}(b_{t-1},f_t,\theta_t,S_t).
% \]

% \smallskip
% \noindent\textbf{(2) Reweight to obtain the next distribution.}
% Let $\mathcal R$ be the set of synergy rules, where each rule is a triple $(T,U,\beta)$ with
% $T\subseteq F$ (trigger), $U\subseteq F$ (target), and $\beta>0$.
% After selecting $f_t$, define the reweighted (unnormalized) probabilities
% \[
% \widehat p_{t+1}(f)
% =
% \bar p_t(f)\cdot
% \prod_{(T,U,\beta)\in\mathcal R:\; f_t\in T,\; f\in U}\beta,
% \qquad f\in F,
% \]
% and then apply feasibility under the updated context $b_t$ and renormalize:
% \[
% p_{t+1}(f)
% =
% \frac{\widehat p_{t+1}(f)}
% {\sum_{g\in F}\widehat p_{t+1}(g)},
% \qquad f\in F.
% \]

\begin{figure}[t]
    \centering
    \includegraphics[width=\linewidth]{images/feature_importance_sql_features_time_rdbms_vs_qiskit.pdf}
    \caption{Importance of various SQL features for SVM and XGBoost trained for optimizing execution time.}
    \label{fig:rdbmsqiskitimportancetime}
\end{figure}
\begin{figure}[t]
    \centering
    \includegraphics[width=\linewidth]{images/feature_importance_sql_features_memory_rdbms_vs_qiskit.pdf}
    \caption{Importance of various SQL features for SVM and RF trained for optimizing memory.}
    \label{fig:rdbmsqiskitimportancemem}
\end{figure}
\section{Additional Results on SQL-only features.}
\label{ssec:sqlf}
Figures~\ref{fig:rdbmsqiskitimportancetime} and~\ref{fig:rdbmsqiskitimportancemem} show that the
dominant SQL signals are structural counts, in particular the number of \texttt{JOIN}s,
predicates, and aggregates. For runtime, the Linear SVM assigns negative coefficients to
\feat{num_where_clauses}, \feat{num_eq_predicates}, and \feat{num_agg_funcs}, suggesting that
queries with heavier filtering and aggregation tend to favor Qiskit Aer. In contrast,
\feat{num_joins} and \feat{num_and_clauses} shift the decision toward the RDBMS backend. For memory,
\feat{num_joins} is the dominant factor, while predicate and aggregation features retain
negative influence, reflecting the benefit of the optimized join algorithms and cost-based planning employed by RDBMS.

% \subsection{Scalability}
% \label{sec:scalability}

% A key motivation behind InferQ is to enable \emph{scalable} data-driven analysis of quantum circuits and simulation workflows. In this section, we evaluate the scalability of InferQ along two complementary dimensions: (i) the cost of \emph{building} the dataset as the number of circuits grows, and (ii) the implications of this cost structure for downstream analytics and learning tasks.

% \subsection{Cumulative Cost of Dataset Construction}

% Figure~\ref{fig:scalability_runtime} shows the cumulative execution time required to process an increasing number of quantum circuits, plotted on a logarithmic scale. We report the cumulative time for three main stages in the InferQ pipeline: circuit generation, feature extraction, and classical simulation. \emph{Several important observations follow.}




% \paragraph{Simulation dominates end-to-end cost.}
% Across the entire range of up to 1{,}000 circuits, the cumulative simulation time grows orders of magnitude faster than both circuit generation and feature extraction. While generation and feature extraction remain within seconds even at scale, simulation reaches the order of $10^3$ seconds. This confirms a central assumption underlying InferQ: \emph{simulation is the primary scalability bottleneck}, whereas data handling and feature computation are comparatively cheap.

% \paragraph{Linear growth in cumulative preprocessing costs.}
% Both circuit generation and feature extraction exhibit smooth, linear growth on the log scale, indicating predictable and well-behaved scaling as more circuits are added. This is particularly important for InferQ, as it suggests that large datasets (tens or hundreds of thousands of circuits) can be constructed incrementally without introducing nonlinear overheads in preprocessing.

% \paragraph{Amortization of expensive simulations.}
% Although simulation is costly, InferQ amortizes this cost by storing simulation outcomes and derived features in a structured database. Once recorded, these results can be reused indefinitely for downstream tasks such as model training, selector construction, and estimator learning, without repeating the simulation. This shifts the scalability challenge from \emph{repeated simulation} to \emph{one-time data acquisition}, a trade-off that strongly favors data-centric workflows.

% \subsection{Implications for Data-Centric Quantum Workflows}

% The scalability behavior observed in Figure~\ref{fig:scalability_runtime} has several implications for the broader use cases presented in Section~\ref{sec:eva}.

% First, it justifies the feasibility of training machine learning models on large-scale simulation data. As shown in the selector and estimator use cases, once simulation results are materialized in InferQ, training and evaluation of complex models (e.g., XGBoost, random forests) become inexpensive relative to the original simulation effort.

% Second, the clear separation between inexpensive static features and expensive dynamic features motivates hybrid approaches. For example, InferQ enables learning-based estimators that predict dynamic properties (such as entropy or sparsity) from static circuit features, thereby avoiding repeated strong simulation for unseen circuits. In this sense, InferQ transforms scalability limits in classical simulation into a supervised learning problem.

% Finally, the use of an RDBMS backend enables horizontal scalability in future deployments. While the experiments in this work were conducted on a single machine, the relational representation of quantum states, operators, and features allows simulation workloads and analytical queries to be distributed across database engines and hardware resources. This opens a path toward scaling InferQ beyond single-node experimentation.

% \subsection{Summary}

% Overall, the scalability evaluation demonstrates that InferQ cleanly isolates the true computational bottleneck—classical simulation—while keeping data generation, feature extraction, and analytics efficient and predictable. By amortizing simulation costs and enabling reuse through structured storage, InferQ provides a scalable foundation for data-driven optimization and analysis of quantum circuit simulation workflows.

\begin{figure}[t]
    \centering
    \includegraphics[width=\linewidth]{images/cumulative_execution_time_generation_vs_feature_extraction.pdf}
    \caption{Cumulative execution time versus number of circuits processed. The figure shows cumulative generation time, and cumulative feature extraction time.}
    \label{fig:scalability_runtime}
\end{figure}

% \subsection{Scalability}
% \label{sec:scalability}

% In Figure~\ref{fig:scalability_runtime}, we report the time changes for generating circuits and extracting features. We observe that generation and feature extraction grow \textbf{linearly} with the increasing number of circuits, even only takes 6.47 seconds to generate 1000 circuits, and extracting features takes 1.9 seconds. 
% Across up to 1{,}000 circuits, cumulative simulation time grows orders of magnitude faster than
% generation and feature extraction. 
% % While generation and feature extraction remain within seconds
% % at this scale, simulation reaches $\sim 10^3$ seconds. This confirms that classical simulation, not
% % data handling, is the primary scalability bottleneck.

% \medskip
% \noindent
% \setlength{\fboxsep}{6pt}
% \fcolorbox{black!15}{black!4}{%
%   \parbox{\dimexpr\linewidth-2\fboxsep-2\fboxrule\relax}{%
% \para{Take-away}
% \sys is scalable.
% }}%

\section{Scalability}
\label{sec:scalability}

Figure~\ref{fig:scalability_runtime} reports how the \sys preprocessing cost grows as we increase the number of generated circuits.
We observe near-linear scaling for both \emph{circuit generation} and \emph{feature extraction} (static, graph, and SQL features): generating
1{,}000 circuits takes only 6.47~seconds, and extracting features takes 1.9~seconds.\footnote{We exclude dynamic features here, as extracting them requires running the simulation, while Section~\ref{ssec:uc2} suggests that they can be predicted using static, graph, and SQL features.}
% In contrast, the cumulative \emph{simulation} time grows orders of magnitude faster than these preprocessing stages.
This tackles Gap 3 in Section~\ref{ssec:limit}. That is, \sys can scale to large circuit corpora: generation and feature extraction contribute negligible overhead.
% , while
% In contrast, simulation remains the dominant cost driver.
% In practice, the workaround is reusing stored results for downstream analysis and learning).

\medskip
  \vspace{0.2cm}
\noindent
\setlength{\fboxsep}{6pt}
\fcolorbox{black!15}{black!4}{%
  \parbox{\dimexpr\linewidth-2\fboxsep-2\fboxrule\relax}{%
\para{Take-away}
\sys scales well in preprocessing: circuit generation and (static/graph/SQL) feature extraction grow linearly and remain inexpensive. This supports building large circuit datasets for database- and learning-based analyses.
}}%



\section{Implementation}
\label{sec:imp}
% Tools and API}

InferQ is supported by two complementary tools:
(1) a \textit{generation and feature-extraction framework} used to construct the dataset, and
(2) a \textit{web-based dataset browser} for interactive exploration and retrieval.

\subsection{InferQ Generation and Feature Extraction Framework}

The parameters in Table~\ref{tab:base_params} are stored in Python dataclass. 

The InferQ framework implements the circuit generation and metric extraction pipeline described in Section~3. It produces combinational quantum circuits and computes their static, graph, SQL, and dynamic features in a reproducible and structured manner.

Each generated circuit is assigned a \emph{content-based hash} that uniquely identifies its structure and parameters. This hash serves as the primary key across the dataset and links the circuit to all extracted feature records.

For every circuit, InferQ stores:
\begin{itemize}
    \item the quantum circuit in a portable serialization format compatible with Qiskit, and
    \item a feature record in JSON format containing all static, graph, and dynamic metrics.
\end{itemize}

Feature extraction is modular and category-based. Static features are derived directly from the circuit representation, graph features are computed from the qubit interaction graph, and dynamic features are obtained from full-state simulation. This separation allows InferQ to be extended with new feature families without recomputing or modifying existing entries.

\subsection{InferQ Web Dataset Viewer}

InferQ is accompanied by a web-based dataset browser inspired by MQTBench\cite{Quetschlich_2023} that enables interactive inspection, filtering, and download of circuits and their features.

The interface presents the dataset as a table where each row corresponds to a circuit identified by its hash. Columns are grouped into the four InferQ feature categories: \textit{static}, \textit{graph}, \textit{SQL}, and \textit{dynamic}.

Users can enable or hide individual feature columns or entire feature groups, apply range-based filters on any visible numeric feature, and explore circuits based on structural, topological, or physical properties.

Each circuit entry includes download links for both the quantum circuit file and its corresponding feature JSON record, allowing users to retrieve exactly the data selected through the interface and use it directly in simulators, learning pipelines, or benchmarking workflows.



\para{Reproducibility and Code} 
InferQ is fully reproducible: given the same code revision, dependency lockfile, and random seed, the framework deterministically regenerates identical circuit structures, parameters, and extracted features. Each circuit is assigned a content-based hash that uniquely identifies its structure and serves as the primary key linking the circuit artifact to all feature records.

\paragraph{Code and environment.}
InferQ is released as open-source software with separate modules for circuit generation, feature extraction, and dataset management. The repository includes a locked Python environment via \texttt{uv.lock} and a declared interpreter version. For exact reproduction, users should report the git commit hash, Python version, and lockfile version.

\paragraph{Cloud-agnostic storage.}
Circuit binaries and feature files are stored in a cloud object store. While Azure Blob Storage is used in the reference deployment, all storage operations are routed through a backend-agnostic interface, allowing any S3-compatible or local filesystem to be used. Each circuit is stored in its own directory
\texttt{Circuits/\{circuit\_hash\}}, containing a serialized \texttt{.qpy} circuit file and a \texttt{.json} metadata file that records generator history and extracted features.


\paragraph{Deterministic circuit generation.}
Circuit construction follows an incremental process in which subcircuits (“generators”) are sampled and appended until a stochastic stopping condition is met. All randomness is controlled by a single global seed, which governs instance-level parameters (e.g., qubit count and depth), generator selection, generator-internal choices, and gate parameters.

\paragraph{Synergy-based generator sampling.}
Generator selection is history-dependent rather than i.i.d. After each generator is chosen, the categorical distribution over the remaining generators is reweighted using predefined \emph{synergies} between generator families. These synergies favor complementary compositions (e.g., state preparation followed by algorithmic cores and optional variational layers), discourage long runs of the same family, and enforce feasibility constraints such as depth and qubit limits. The updated distribution is then renormalized before sampling the next generator. Because these updates are deterministic functions of the previous generator sequence and the global seed, the entire generator sequence is exactly reproducible.

\paragraph{Feature reproducibility.}
For each circuit, InferQ extracts static, graph-based, and dynamic (simulation-based) features. Feature extractors are modular and deterministic given the same simulator backend, numeric precision, and hardware configuration. For fair comparison, users should report simulator version and floating-point settings when dynamic features are used.


\subsection{Qiskit Aer Data Representation Selector}
\label{ssec:selector_RESULTS}
Most quantum applications will be run using a simulator framework like IBM Qiskit \cite{IBMQiskit} with Qiskit Aer\cite{QiskitAerDocs} as a simulation backend. Qiskit Aer offers various simulation backends corresponding to different quantum data representations. Each performs better or worse in terms of time and memory depending on the circuit to be simulated. This leads to the research question: \textit{what is the optimal Qiskit Aer data representation for classical simulation?} 

InferQ can be used to train a selector for the Qiskit data representations for optimal simulation. In our experiment we trained various ML models  to select the best representation within the classes of ["statevector", "density\_matrix", "extended\_stabilizer", "matrix\_product\_state"] provided by Qiskit Aer. We recorded the execution time for 93{,}567 entries in InferQ on our machine from Section~\ref{sec:impl}.\footnote{Qiskit Selector training data (Parquet) available at \url{https://github.com/InfiniData-Lab/InferQ/blob/main/analysis/training_data/selector_training_data.parquet}} For this experiment, we used the framework from Qymera \cite{Littau_2025}.\footnote{We thank the authors for sharing their code and helping us set up the experiments.}
 

By defining the best representation to be the one taking lowest execution time, we could train various multiclass classifier models to select the best method. 

Training models of complexities and  evaluating them on a test set of size 20\% outperformed Qiskit Aer's own ``automatic'' data representation selector by a margin of up to 28\%.

The accuracy in table  \ref{tab:selectoraccuracy} is the percentage of cases in the test set that the predicted best method matched the actual best method. Upon inspection of the features, it became clear why Naive Bayes performs weakly. There is multimodality within feature distributions rather than Gaussian structures.  In addition, the classes are highly imbalanced, as shown in Figure \ref{fig:selectortestset}.

\begin{table}[tb]
\centering
\begin{minipage}{\columnwidth}

\caption{Test-set accuracy of InferQ-trained selectors compared with Qiskit Aer's built-in ``automatic'' method selector.}
\label{tab:selectoraccuracy}

\centering
\resizebox{\columnwidth}{!}{%
\begin{tabular}{cccccccc}
\toprule
 & Qiskit ``Automatic'' & \multicolumn{6}{c}{Trained on InferQ} \\
Metric & Accuracy & \textbf{XGBoost} & Random Forest & Decision Tree & KNN & Logistic Regression & Naive Bayes \\
\midrule
Accuracy & 36\% & \textbf{64\%} & 64\% & 63\% & 61\% & 61\% & 55\% \\
\bottomrule
\end{tabular}
}

\vspace{0.6em}

\centering
\includegraphics[width=0.7\linewidth]{images/selector-distribution.png}
\captionof{figure}{Bar chart displaying the imbalance between the classes being the best simulation method for minimum execution time.}
\label{fig:selectortestset}

\end{minipage}
\end{table}

\begin{figure*}[tbh]
  \centering
  \includegraphics[width=\linewidth]{images/finalinferQschema.pdf}
  \caption{Schema for storing extracted circuit features. 
  \sys circuit data model schema. Each circuit is stored with a unique
  hash (primary key) and associated feature classes. \sys entries are centered
  around circuits in the \texttt{InferQ\_Circuits} table, which reference
  static, graph, SQL and dynamic features, as well as storage metadata in Azure DB. 
      }
  \label{fig:schema}
\end{figure*}

Despite the K-Nearest Neighbour (KNN)'s relatively high accuracy, it is clear from 
Figure \ref{fig:selectorconfusion} 
that it is unable to predict the smaller classes well, as 42\% of the ``density\_matrix'' and 34\% ``extended\_stabilizer'' methods are being classified as ``matrix\_product\_state''. The Decision Tree is much better overall for all the classes, except it is overcompensating for the minority classes and hence has low accuracy for the dominant ``matrix\_product\_state''. Logistic regression has lowest accuracy so we are left with tree models and linear SVM.

\begin{figure*}[tb]
    \centering
    \includegraphics[width=\linewidth]{images/selector-confusion_1.png}
    \caption{Confusion matrices, normalised over the true classes, for each model's performance on the test set.
    SV = \texttt{statevector}, DM = \texttt{density\_matrix}, MPS = \texttt{matrix\_product\_state}, ES = \texttt{extended\_stabilizer}.
    Panels are ordered by decreasing accuracy and share a common colour scale.}
    \label{fig:selectorconfusion}
\end{figure*}

In order to break this tie, we can look at the distribution in which the simulation methods come 2nd, 3rd and 4th place. From Table \ref{tab:method_ranks} it becomes clear that ``matrix\_product\_state'', ``density\_matrix'' and "statevector" dominate in 2nd position too. Hence, we can choose XGBoost which performs the best for these 3 dominant classes as can be seen from confusion plots of Figure 33. Although the accuracy of XGBoost is 64.4\%, it still outperforms Qiskit Aer's own "automatic" method which is only able to select correctly in 35.9\% of the cases. Even simpler models like logistic regression and linear SVMs are able to leverage InferQ features to perform relatively well against the current Qiskit Aer "automatic" selector. 

\begin{table}[h]
\caption{Number of entries where each simulation method is ranked 1st, 2nd, 3rd, and 4th by the XGBoost-based selector, ordered by predicted execution time.}
\label{tab:method_ranks}
\centering
\resizebox{\columnwidth}{!}{%
\begin{tabular}{lcccc}
\toprule
 & \multicolumn{4}{c}{Rank Position} \\
\cmidrule(lr){2-5}
Method & 1 & 2 & 3 & 4 \\
\midrule
\texttt{statevector}            & 30802 & 34743 & 25131 & 2891 \\
\texttt{density\_matrix}        & 17704 & 20340 & 7625  & 948  \\
\texttt{matrix\_product\_state} & 42288 & 35959 & 14537 & 783  \\
\texttt{extended\_stabilizer}   & 2773  & 2525  & 46274 & 41995 \\
\bottomrule
\end{tabular}%
}
\end{table}


We thus show that InferQ can be used to optimize Qiskit Aer Workloads by building a selector, demonstrating that a complex model like XGBoost surpasses its simpler competitors in predictions amongst the highly imbalanced classes of representations. By outperforming Qiskit Aer's own "automatic" selector with all models experimented with, it validates our data-centric approach to this problem.

Moreover, the experiment  showcases that data management techniques can be used to make classical simulation of quantum computation more efficient without incorporating the domain specific techniques of quantum science. It also completes a pipeline to solve the mutli-class problem of which classical simulation method is optimal given a circuit incorporating the novel approach of RDBMS.

\section{InferQ Feature Schema}
\label{sec:inferqschemafull}
Figure~\ref{fig:schema} illustrates the schema of how \sys stores the extracted circuit features. 







\end{appendices}